% =====================================================================
%  RAPPORT DE PROJET DE FIN D'ETUDES
%  XPARK GMAO - plateforme de gestion de flotte et de transport
% =====================================================================
%  Compilation :  pdflatex main  (deux fois, pour la table des matieres,
%                 la table des figures et la pagination n/N)
%  Ou simplement :  compiler.bat
% =====================================================================
\documentclass[12pt,a4paper,twoside,openright]{report}

% =====================================================================
%  Paquets
% =====================================================================
%  Compilateur retenu : pdfLaTeX. Times est fourni par mathptmx, sans
%  dependance a une police du systeme : le document se compile a
%  l'identique sur n'importe quelle installation TeX.
% =====================================================================

\usepackage[utf8]{inputenc}
\usepackage[T1]{fontenc}
\usepackage[french]{babel}
\usepackage{mathptmx}                 % Times New Roman pour le corps
\usepackage[scaled=0.90]{helvet}      % Helvetica pour les titres
\usepackage{courier}

\usepackage[a4paper,margin=2.5cm,headheight=15pt]{geometry}
\usepackage{setspace}
\usepackage{graphicx}
\usepackage{xcolor}
\usepackage{float}
\usepackage{pdflscape}                % diagrammes larges en paysage
\usepackage{booktabs}
\usepackage{longtable}
\usepackage{array}
\usepackage{multirow}
\usepackage{tabularx}
\usepackage[font=small,labelfont=bf,labelsep=period]{caption}
\usepackage{subcaption}
\usepackage{enumitem}
\usepackage{fancyhdr}
\usepackage{lastpage}                 % pagination sous la forme n/N
\usepackage{titlesec}
\usepackage{tocloft}
\usepackage[hidelinks]{hyperref}
\usepackage{url}
\usepackage{tcolorbox}
\usepackage{tikz}                     % pose des champs sur la couverture

\tcbuselibrary{skins,breakable}
\usetikzlibrary{calc}

% =====================================================================
%  Mise en forme imposee par « Rapport Stage Pedagogique.pdf »
% =====================================================================
%  Police Times 12 pt, marges 2,5 cm, interligne 1,15, justification,
%  alinea de premiere ligne 0,50 cm, pagination en bas a droite sous la
%  forme n/N. Chaque valeur ci-dessous est reprise du guide.
% =====================================================================

\setstretch{1.15}
\setlength{\parindent}{0.5cm}
\setlength{\parskip}{0.15\baselineskip}

% ---------------------------------------------------------------- Titres
\definecolor{bleuTitre}{HTML}{1F3864}
\definecolor{grisFilet}{HTML}{9AA5B1}
\definecolor{grisFond}{HTML}{F2F4F7}

\titleformat{\chapter}[display]
  {\normalfont\sffamily\bfseries\color{bleuTitre}}
  {\filright\large\MakeUppercase{\chaptertitlename}\ \Huge\thechapter}
  {8pt}{\huge\filright}
\titlespacing*{\chapter}{0pt}{-20pt}{28pt}

\titleformat{\section}{\normalfont\sffamily\large\bfseries\color{bleuTitre}}
  {\thesection}{0.6em}{}
\titleformat{\subsection}{\normalfont\sffamily\normalsize\bfseries\color{bleuTitre}}
  {\thesubsection}{0.6em}{}
\titleformat{\subsubsection}{\normalfont\sffamily\small\bfseries\color{bleuTitre}}
  {\thesubsubsection}{0.6em}{}

\titlespacing*{\section}      {0pt}{14pt}{6pt}
\titlespacing*{\subsection}   {0pt}{11pt}{4pt}
\titlespacing*{\subsubsection}{0pt}{9pt}{3pt}

\setcounter{secnumdepth}{3}
\setcounter{tocdepth}{2}

% ---------------------------------------------------------------- En-tetes
\pagestyle{fancy}
\fancyhf{}
\renewcommand{\headrulewidth}{0.4pt}
\renewcommand{\footrulewidth}{0pt}
\fancyhead[L]{\small\sffamily\nouppercase{\leftmark}}
\fancyfoot[R]{\small\sffamily\thepage/\pageref{LastPage}}

\fancypagestyle{plain}{%
  \fancyhf{}%
  \renewcommand{\headrulewidth}{0pt}%
  \fancyfoot[R]{\small\sffamily\thepage/\pageref{LastPage}}%
}

\renewcommand{\chaptermark}[1]{\markboth{\thechapter.\ #1}{}}

% ---------------------------------------------------------------- Legendes
\captionsetup[figure]{name=Figure}
\captionsetup[table]{name=Tableau}
\renewcommand{\cftfigpresnum}{Figure~}
\renewcommand{\cfttabpresnum}{Tableau~}
\setlength{\cftfignumwidth}{4.2em}
\setlength{\cfttabnumwidth}{4.6em}

% Les tableaux portent leur legende au-dessus, les figures en dessous.
\captionsetup[table]{position=above,skip=4pt}
\captionsetup[figure]{position=below,skip=6pt}

% ---------------------------------------------------- Pages liminaires
%  POURQUOI DES PAGES BLANCHES APPARAISSENT
%  La classe est chargee en « twoside, openright » : chaque nouvelle
%  piece commence sur une page de DROITE, comme il est d'usage pour un
%  document imprime en recto-verso. Quand la piece precedente se termine
%  elle-meme sur une page de droite, LaTeX intercale une page blanche.
%  C'est voulu, et c'est ce que demande le guide, qui recommande
%  l'impression recto-verso.
%
%  Les CHAPITRES conservent ce comportement : commencer un chapitre au
%  dos d'un autre serait fautif.
%
%  Les pages liminaires, elles, n'ont pas cette exigence. Le reglage
%  ci-dessous permet de choisir :
%      \clearpage        -> pas de page blanche entre elles  (defaut)
%      \cleardoublepage  -> chacune sur une page de droite
\newcommand{\sautLiminaire}{\clearpage}

% ---------------------------------------------------------------- Divers
\renewcommand{\arraystretch}{1.25}
\setlist{itemsep=2pt,parsep=0pt,topsep=4pt}
\clubpenalty=10000      % pas de ligne isolee en bas de page
\widowpenalty=10000     % pas de ligne isolee en haut de page
\raggedbottom

% =====================================================================
%  Macros de figures
% =====================================================================
%  Toutes les figures du rapport passent par ces macros : la mise en
%  forme (centrage, largeur, legende, espacement) reste ainsi identique
%  d'un bout a l'autre du document.
%
%  Le guide impose que chaque figure porte une legende ET soit citee
%  dans le texte. Chaque macro prend donc un label, et le texte doit
%  contenir un \ref correspondant.
% =====================================================================

\graphicspath{{figures/}{figures/diagrammes/}{figures/logos/}%
              {figures/captures/}{figures/planning/}}

% ---------------------------------------------------------------------
%  \figReelle{fichier}{legende}{label}{largeur}
%  Figure ordinaire. La largeur s'exprime en fraction de \textwidth.
% ---------------------------------------------------------------------
%  Placement « htbp » et non « H » : un diagramme agrandi ne tient
%  souvent pas dans le bas de page ou il tombe. Ancre de force, il y
%  laisse un vide et se reporte a la page suivante ; laisse flottant, il
%  se pose en tete de page et le texte comble l'intervalle. Le gain
%  mesure sur ce document est de sept pages, sans perte de proximite
%  puisque chaque diagramme est cite dans le texte par son numero.
\newcommand{\figReelle}[4]{%
  \begin{figure}[htbp]
    \centering
    \includegraphics[width=#4\textwidth,keepaspectratio]{#1}
    \caption{#2}
    \label{#3}
  \end{figure}%
}

% ---------------------------------------------------------------------
%  \capture{fichier}{legende}{label}
%  Capture d'ecran : large, encadree d'un filet fin pour la detacher du
%  fond de page. Le regroupement de quatre ou cinq captures sous une
%  seule figure a ete abandonne : reduites a moins de la moitie de la
%  largeur du texte, elles devenaient illisibles a l'impression.
%
%  LARGEUR UNIQUE. Les trente-sept captures du rapport ont ete prises
%  dans la meme fenetre : elles partagent le meme rapport, 2,07. Une
%  largeur unique suffit donc, et c'est elle qui donne a une page
%  portant deux captures son aplomb. Toute exception se verrait.
% ---------------------------------------------------------------------
\newcommand{\largeurCapture}{0.90}

\newcommand{\capture}[3]{%
  \begin{figure}[H]
    \centering
    \fbox{\includegraphics[width=\largeurCapture\textwidth,keepaspectratio]{#1}}
    \caption{#2}
    \label{#3}
  \end{figure}%
}

%  Conservee pour compatibilite : elle composait autrefois plus etroit.
%  Elle ne le fait plus. La seule capture qui l'employait, celle de
%  l'application chauffeur, mesure 1887 x 912 pixels comme les autres :
%  la reduire a 0,60 la rendait simplement plus petite que ses voisines,
%  sans raison.
\newcommand{\captureEtroite}[3]{\capture{#1}{#2}{#3}}

% ---------------------------------------------------------------------
%  \figDeux / \figTrois : plusieurs images sous UNE seule legende.
%  Reservees aux DIAGRAMMES qui se lisent ensemble. Ne plus les employer
%  pour des captures d'ecran.
% ---------------------------------------------------------------------
\newcommand{\figDeux}[5]{%
  \begin{figure}[H]
    \centering
    \includegraphics[width=#5\textwidth,keepaspectratio]{#1}\\[4pt]
    \includegraphics[width=#5\textwidth,keepaspectratio]{#2}
    \caption{#3}
    \label{#4}
  \end{figure}%
}

\newcommand{\figTrois}[6]{%
  \begin{figure}[H]
    \centering
    \includegraphics[width=#6\textwidth,keepaspectratio]{#1}\\[3pt]
    \includegraphics[width=#6\textwidth,keepaspectratio]{#2}\\[3pt]
    \includegraphics[width=#6\textwidth,keepaspectratio]{#3}
    \caption{#4}
    \label{#5}
  \end{figure}%
}

% ---------------------------------------------------------------------
%  \figPaysage{fichier}{legende}{label}
%  Diagramme trop large pour la justification : page en paysage.
% ---------------------------------------------------------------------
\newcommand{\figPaysage}[3]{%
  \begin{landscape}
    \begin{figure}[H]
      \centering
      \includegraphics[width=0.95\linewidth,height=0.82\textheight,%
                       keepaspectratio]{#1}
      \caption{#2}
      \label{#3}
    \end{figure}
  \end{landscape}%
}

% =====================================================================
%  Emplacements reserves
% =====================================================================
%  Une figure encore absente occupe deja sa place, avec sa VRAIE
%  legende et son VRAI label : elle est donc numerotee, citee dans le
%  texte et listee dans la table des figures. Le jour ou l'image
%  existe, seule la commande change : \figAttendue devient \figReelle.
% =====================================================================
\definecolor{cadreAttente}{HTML}{B0562A}
\definecolor{fondAttente}{HTML}{FDF3EC}

\newcommand{\figAttendue}[3]{%
  \begin{figure}[H]
    \centering
    \begin{tcolorbox}[width=0.86\textwidth,height=5.4cm,
      colback=fondAttente,colframe=cadreAttente,boxrule=1pt,
      borderline={0.5pt}{3pt}{cadreAttente,dashed},
      halign=center,valign=center,sharp corners]
      {\sffamily\bfseries\color{cadreAttente}FIGURE A FOURNIR}\\[6pt]
      {\sffamily\small Fichier attendu :}\\[2pt]
      {\ttfamily\small #1}
    \end{tcolorbox}
    \caption{#2}
    \label{#3}
  \end{figure}%
}

%  Variante pour les captures d'ecran : indique en plus l'ecran a
%  photographier, pour que la reprise soit immediate.
\newcommand{\captureAttendue}[4]{%
  \begin{figure}[H]
    \centering
    \begin{tcolorbox}[width=0.86\textwidth,height=5.4cm,
      colback=fondAttente,colframe=cadreAttente,boxrule=1pt,
      borderline={0.5pt}{3pt}{cadreAttente,dashed},
      halign=center,valign=center,sharp corners]
      {\sffamily\bfseries\color{cadreAttente}CAPTURE A FOURNIR}\\[6pt]
      {\sffamily\small Ecran :} {\sffamily\itshape #4}\\[4pt]
      {\sffamily\small Fichier attendu :}\\[2pt]
      {\ttfamily\small #1}
    \end{tcolorbox}
    \caption{#2}
    \label{#3}
  \end{figure}%
}

% =====================================================================
%  Encadres de contenu
% =====================================================================
\newtcolorbox{definitionbox}[1]{%
  breakable, colback=grisFond, colframe=grisFilet, boxrule=0.6pt,
  left=8pt, right=8pt, top=6pt, bottom=6pt, sharp corners,
  title={\sffamily\bfseries #1}, coltitle=bleuTitre,
  colbacktitle=grisFond, titlerule=0pt}

% Mise en valeur imposee par le guide : les exemples en italique, les
% definitions encadrees, les elements essentiels en gras.
\newcommand{\exemple}[1]{\textit{#1}}
\newcommand{\essentiel}[1]{\textbf{#1}}

% =====================================================================
%  Tableau de backlog de sprint
% =====================================================================
\newenvironment{backlog}[1]%
  {\begin{table}[H]\centering\caption{#1}%
   \small\begin{tabularx}{\textwidth}{@{}c>{\raggedright\arraybackslash}X c c@{}}
   \toprule
   \textbf{N\textsuperscript{o}} & \textbf{Recit utilisateur} &
   \textbf{Priorite} & \textbf{Estimation} \\ \midrule}%
  {\bottomrule\end{tabularx}\end{table}}

% =====================================================================
%  DONNEES VARIABLES DU RAPPORT
% =====================================================================
%  Tout ce qui depend de vous et de l'entreprise est regroupe ici.
%  C'est le SEUL fichier a modifier pour personnaliser le document :
%  page de garde, formulaire de depot, remerciements et resume s'en
%  servent tous.
%
%  Les valeurs marquees A COMPLETER apparaissent en rouge dans le PDF,
%  ce qui rend impossible de les oublier avant impression.
% =====================================================================

\newcommand{\aCompleter}[1]{\textcolor{red}{\textbf{[#1]}}}

% ---------------------------------------------------------------- Etudiant
\newcommand{\etudiantPrenom}{Saber}
\newcommand{\etudiantNom}{Barhoumi}
\newcommand{\etudiantComplet}{\etudiantPrenom{} \etudiantNom}

% ---------------------------------------------------------------- Formation
\newcommand{\anneeUniversitaire}{2025 - 2026}
\newcommand{\filiere}{INFORMATIQUE}
\newcommand{\etablissement}{ESPRIT}
\newcommand{\etablissementComplet}{Ecole Superieure Privee d'Ingenierie
  et de Technologies}
\newcommand{\etablissementLogo}{logo-esprit.png}

% ---------------------------------------------------------------- Sujet
%  Titre officiel, tel qu'il doit figurer sur la couverture. Il porte
%  ses accents, contrairement au corps du rapport qui est ecrit sans :
%  un titre officiel se reproduit a la lettre.
%
%  Deux formes. La premiere, d'un seul tenant, sert aux proprietes du
%  fichier PDF et au resume. La seconde porte ses coupures de ligne :
%  sur une couverture, un titre ne se coupe pas au milieu d'un mot, et
%  les parametres de cesure ne traversent pas un noeud TikZ centre.
\newcommand{\sujetRapport}{Conception et développement d'une application
  GMAO pour la gestion de flotte --- Cas de XPARK}

\newcommand{\sujetRapportCouverture}{%
  Conception et développement d'une\\
  application GMAO pour la gestion de flotte\\
  --- Cas de XPARK ---}
\newcommand{\sujetCourt}{XPARK GMAO}

% ---------------------------------------------------------------- Entreprise
\newcommand{\entreprise}{SIPRO}
\newcommand{\entrepriseSecteur}{les technologies de l'information}
\newcommand{\entrepriseVille}{Sfax}
\newcommand{\entrepriseCodePostal}{3000}
\newcommand{\entrepriseLogo}{logo-sipro.png}

% ---------------------------------------------------------------- Encadrants
\newcommand{\encadrantAcademique}{M. Karray Gargouri}
\newcommand{\encadrantEntreprise}{M. Salah Tabib}

% ---------------------------------------------------------------- Stage
\newcommand{\stageDebut}{15 fevrier 2026}
\newcommand{\stageFin}{15 aout 2026}
\newcommand{\stageDuree}{six}

% ---------------------------------------------------------------- Depot
\newcommand{\depotGitHub}{https://github.com/saber-barhoumi/Xpark-gmao}


\hypersetup{
  pdftitle={\sujetRapport},
  pdfauthor={\etudiantComplet},
  pdfsubject={Rapport de projet de fin d'etudes},
  pdfkeywords={microservices, gestion de flotte, transport, Scrum, .NET, React}
}

\begin{document}

% =====================================================================
%  PAGES LIMINAIRES  (numerotation romaine)
% =====================================================================
\pagenumbering{roman}

% =====================================================================
%  Page de garde, reprise du gabarit « Page Garde Rapport Stage(Fr) »
% =====================================================================
\begin{titlepage}
\thispagestyle{empty}
\centering
\sffamily

\begin{minipage}[c]{0.30\textwidth}
  \raggedright
  \includegraphics[height=2.1cm,keepaspectratio]{\etablissementLogo}
\end{minipage}
\hfill
\begin{minipage}[c]{0.34\textwidth}
  \centering
  {\large\bfseries \anneeUniversitaire}
\end{minipage}
\hfill
\begin{minipage}[c]{0.30\textwidth}
  \raggedleft
  \includegraphics[height=2.1cm,keepaspectratio]{\entrepriseLogo}
\end{minipage}

\vspace{1.1cm}
{\Large\bfseries\color{bleuTitre} \filiere\par}

\vspace{0.45cm}
{\color{grisFilet}\rule{\textwidth}{1.2pt}}

\vspace{1.5cm}
{\bfseries\large RAPPORT DE PROJET DE FIN D'ETUDES\par}

\vspace{1.1cm}
{\fontsize{22}{28}\selectfont\bfseries\color{bleuTitre}
 \sujetRapport\par}

\vspace{0.7cm}
{\Large\itshape \sujetCourt\par}

\vspace{1.3cm}
{\color{grisFilet}\rule{0.6\textwidth}{0.8pt}}

\vspace{1.3cm}
\begin{minipage}{0.9\textwidth}
  \raggedright\large
  \textbf{Realise par :} \etudiantComplet\\[10pt]
  \textbf{Encadrant academique :} \encadrantAcademique\\[6pt]
  \textbf{Encadrant entreprise :} \encadrantEntreprise\\[10pt]
  \textbf{Organisme d'accueil :} \entreprise{} -- \entrepriseVille
\end{minipage}

\vfill

{\small \etablissementComplet\par}
\vspace{0.3cm}
{\small Annee universitaire \anneeUniversitaire}

\end{titlepage}

% =====================================================================
%  Formulaire de validation de depot
%  Le gabarit precise : « Ce formulaire doit etre rempli, signe et
%  scanne » et « introduit apres la page de garde ».
% =====================================================================
\cleardoublepage
\thispagestyle{empty}
\sffamily
\vspace*{1.5cm}

\begin{center}
  {\Large\bfseries\color{bleuTitre}
   Validation du depot du rapport de projet de fin d'etudes\par}
  \vspace{4pt}
  {\large\itshape Report submission validation\par}
\end{center}

\vspace{1.2cm}

\noindent
Je valide le depot du rapport de projet de fin d'etudes relatif a
l'etudiant nomme ci-dessous.\\
\textit{I validate the submission of the student's report:}

\vspace{0.8cm}
\noindent
\textbf{Nom et prenom} \textit{/ Name and surname} :
\etudiantComplet

\vspace{1.6cm}
{\color{grisFilet}\rule{\textwidth}{0.6pt}}
\vspace{0.8cm}

\noindent
\begin{minipage}[t]{0.48\textwidth}
  \textbf{Encadrant entreprise}\\
  \textit{Business site supervisor}\\[10pt]
  Nom et prenom : \encadrantEntreprise\\[1.9cm]
  Cachet et signature\\
  \textit{Stamp and signature}\\[0.4cm]
  {\color{grisFilet}\rule{0.9\linewidth}{0.5pt}}
\end{minipage}
\hfill
\begin{minipage}[t]{0.48\textwidth}
  \textbf{Encadrant academique}\\
  \textit{Academic supervisor}\\[10pt]
  Nom et prenom : \encadrantAcademique\\[1.9cm]
  Signature\\
  \textit{Signature}\\[0.4cm]
  {\color{grisFilet}\rule{0.9\linewidth}{0.5pt}}
\end{minipage}

\vfill
\begin{center}
  \small\itshape
  Ce formulaire doit etre rempli, signe et scanne.\\
  This form must be completed, signed and scanned.
\end{center}
\vspace{1cm}

\rmfamily

% =====================================================================
%  Dedicaces
% =====================================================================
%  A REECRIRE : ce texte est une trame, remplacez-le par le votre.
%  Le guide ne l'impose pas mais le rapport de reference en comporte une.
% =====================================================================
\cleardoublepage
\thispagestyle{empty}
\vspace*{2.5cm}

\begin{center}
  {\Huge\sffamily\bfseries\color{bleuTitre} Dedicaces\par}
\end{center}

\vspace{2cm}

\begin{center}
\begin{minipage}{0.78\textwidth}
\itshape\setstretch{1.4}

Du plus profond de mon coeur, je dedie ce travail

\vspace{1cm}
\centering
\textbf{A mes parents,}\\[6pt]
\raggedright
pour leur soutien sans faille et les sacrifices consentis tout au long
de mes etudes. Rien de ce qui suit n'aurait ete possible sans eux.

\vspace{1cm}
\centering
\textbf{A ma famille et a mes proches,}\\[6pt]
\raggedright
pour leurs encouragements constants et leur patience pendant les longues
periodes de travail.

\vspace{1cm}
\centering
\textbf{A mes enseignants,}\\[6pt]
\raggedright
qui m'ont transmis le gout du travail rigoureux et les fondements sur
lesquels repose ce projet.

\vspace{1cm}
\centering
\textbf{A mes amis,}\\[6pt]
\raggedright
pour leur presence et leur bonne humeur, qui ont rendu ce parcours plus
leger.

\end{minipage}
\end{center}

\vfill
\begin{flushright}
  \itshape \etudiantComplet
\end{flushright}
\vspace{1cm}

% =====================================================================
%  Remerciements
% =====================================================================
%  Le guide les rend obligatoires : « Il est d'usage de mettre egalement
%  une page pour remercier tous ceux envers qui vous etes redevables ».
%  A RELIRE ET PERSONNALISER une fois les noms connus.
% =====================================================================
\cleardoublepage
\phantomsection
\addcontentsline{toc}{chapter}{Remerciements}
\thispagestyle{empty}
\vspace*{1.5cm}

\begin{center}
  {\Huge\sffamily\bfseries\color{bleuTitre} Remerciements\par}
\end{center}

\vspace{1.5cm}

\setstretch{1.3}

Au terme de ce projet de fin d'etudes, je tiens a exprimer ma
reconnaissance a toutes les personnes qui ont contribue, de pres ou de
loin, a son aboutissement.

\vspace{0.6cm}

J'adresse mes remerciements les plus sinceres a
\encadrantAcademique, mon encadrant academique, pour la qualite de son
suivi, la pertinence de ses remarques et la disponibilite dont il a fait
preuve tout au long de ce travail. Ses conseils methodologiques ont
oriente la demarche suivie et ont contribue a la rigueur de ce rapport.

\vspace{0.6cm}

Je remercie egalement \encadrantEntreprise, mon encadrant au sein de
\entreprise, pour la confiance accordee, pour l'autonomie laissee dans
les choix techniques et pour l'accompagnement quotidien qui a permis de
mener ce projet a son terme.

\vspace{0.6cm}

Ma gratitude va aussi a l'ensemble de l'equipe de \entreprise, dont la
disponibilite et les retours d'experience metier ont ete determinants
pour comprendre les besoins reels de la gestion de flotte et du
transport de marchandises.

\vspace{0.6cm}

Je remercie le corps enseignant de \etablissement{} pour la formation
recue durant ces annees d'etudes, ainsi que les membres du jury qui me
font l'honneur d'evaluer ce travail.

\vspace{0.6cm}

Enfin, mes remerciements les plus chaleureux vont a ma famille et a mes
proches pour leur soutien constant.

\setstretch{1.15}
\vspace{1cm}


\cleardoublepage
\phantomsection
\addcontentsline{toc}{chapter}{Table des matieres}
\tableofcontents

\cleardoublepage
\phantomsection
\addcontentsline{toc}{chapter}{Table des figures}
\listoffigures

\cleardoublepage
\phantomsection
\addcontentsline{toc}{chapter}{Liste des tableaux}
\listoftables

% =====================================================================
%  Liste des acronymes
% =====================================================================
%  Volontairement tenue a la main plutot qu'avec le paquet glossaries :
%  cela evite l'etape makeglossaries et garde la compilation a un simple
%  pdflatex lance deux fois.
% =====================================================================
\sautLiminaire
\phantomsection
\addcontentsline{toc}{chapter}{Liste des acronymes}
\chapter*{Liste des acronymes}
\markboth{Liste des acronymes}{}

\begin{longtable}{@{}p{2.6cm}p{11.5cm}@{}}
\toprule
\textbf{Sigle} & \textbf{Signification} \\ \midrule
\endfirsthead
\toprule
\textbf{Sigle} & \textbf{Signification} \\ \midrule
\endhead

AMQP    & \textit{Advanced Message Queuing Protocol}, protocole d'echange de messages \\
API     & \textit{Application Programming Interface}, interface de programmation \\
CORS    & \textit{Cross-Origin Resource Sharing}, partage de ressources entre origines \\
CRUD    & \textit{Create, Read, Update, Delete}, les quatre operations de base sur une donnee \\
DTO     & \textit{Data Transfer Object}, objet de transfert de donnees \\
EF Core & \textit{Entity Framework Core}, couche de correspondance objet-relationnel de .NET \\
GMAO    & Gestion de la maintenance assistee par ordinateur \\
GPS     & \textit{Global Positioning System}, systeme de positionnement par satellite \\
HTTP    & \textit{HyperText Transfer Protocol} \\
IA      & Intelligence artificielle \\
IAM     & \textit{Identity and Access Management}, gestion des identites et des acces \\
IBOA    & \textit{Intelligent Business Owner Assistant}, assistant de veille metier \\
JSON    & \textit{JavaScript Object Notation}, format d'echange de donnees \\
JWT     & \textit{JSON Web Token}, jeton d'authentification signe \\
MFA     & \textit{Multi-Factor Authentication}, authentification a plusieurs facteurs \\
OT      & Ordre de travail \\
PFE     & Projet de fin d'etudes \\
REST    & \textit{Representational State Transfer}, style d'architecture d'API \\
RAG     & \textit{Retrieval-Augmented Generation}, generation augmentee par recherche \\
SGBD    & Systeme de gestion de base de donnees \\
SPA     & \textit{Single Page Application}, application web monopage \\
SQL     & \textit{Structured Query Language}, langage de requetes \\
TVA     & Taxe sur la valeur ajoutee \\
UML     & \textit{Unified Modeling Language}, langage de modelisation unifie \\
URL     & \textit{Uniform Resource Locator}, adresse d'une ressource web \\
VIN     & \textit{Vehicle Identification Number}, numero de chassis \\
XGBoost & \textit{eXtreme Gradient Boosting}, algorithme d'apprentissage supervise \\
\bottomrule
\end{longtable}


% =====================================================================
%  CORPS DU RAPPORT  (numerotation arabe)
% =====================================================================
\cleardoublepage
\pagenumbering{arabic}
\setcounter{page}{1}

% =====================================================================
%  Introduction generale
% =====================================================================
%  Le guide precise : « pose le probleme a resoudre sans evocation de
%  resultats », puis « le contenu de chaque chapitre est annonce
%  brievement ».
% =====================================================================
\cleardoublepage
\phantomsection
\addcontentsline{toc}{chapter}{Introduction generale}
\chapter*{Introduction generale}
\markboth{Introduction generale}{}

Le transport routier de marchandises occupe une place structurante dans
l'economie tunisienne. Il relie les zones industrielles du littoral aux
bassins de production de l'interieur, dessert les ports et assure la
distribution vers les reseaux commerciaux. Les entreprises du secteur
exploitent des flottes de plusieurs dizaines de vehicules, mobilisent des
chauffeurs sur des trajets longs et repondent a des clients qui attendent
des delais fermes et une facturation exacte.

Cette activite repose sur un equilibre fragile. Un vehicule immobilise
par une panne non anticipee, un document administratif expire decouvert
au moment d'un controle, un devis etabli sans reference tarifaire fiable
ou une facture oubliee : chacun de ces incidents a un cout direct et
degrade la relation commerciale. Les entreprises de taille moyenne
gerent pourtant ces risques avec des outils qui n'ont pas ete concus
pour cela, principalement des tableurs et des echanges informels.

\vspace{0.4cm}
\noindent
\essentiel{La question a laquelle ce projet de fin d'etudes entend
repondre est la suivante : comment doter une entreprise de transport de
marchandises d'un systeme d'information unifie qui centralise
l'exploitation de sa flotte, industrialise sa maintenance, ouvre un canal
numerique a ses clients et mette l'aide a la decision au service de ses
choix quotidiens ?}
\vspace{0.4cm}

Repondre a cette question suppose de traiter simultanement plusieurs
domaines metier qui, dans les pratiques observees, sont geres
separement : les identites et les habilitations, le referentiel des
vehicules et des equipements, le cycle commercial de la demande a la
facture, la maintenance curative et preventive, la communication interne
et la relation client. L'enjeu n'est pas seulement de les informatiser un
a un, mais de les faire cooperer sans les rendre indissociables : une
evolution de la facturation ne doit pas immobiliser la gestion de la
flotte.

Le travail presente ici a consiste a concevoir et a developper
\sujetCourt, une plateforme construite selon une architecture
microservices. Elle reunit douze services autonomes, trois applications
destinees a des publics distincts et une couche d'aide a la decision
fondee sur l'apprentissage automatique et sur la recherche documentaire
augmentee. La demarche suivie est agile, structuree par le cadre Scrum
et decoupee en six sprints correspondant chacun a un domaine metier
coherent.

\section*{Organisation du rapport}
\addcontentsline{toc}{section}{Organisation du rapport}

Ce rapport suit la progression reelle du projet.

Le \essentiel{premier chapitre} situe le travail dans son contexte. Il
presente l'organisme d'accueil, formule la problematique, examine les
pratiques et les solutions existantes, expose leurs limites, decrit la
solution proposee et justifie le choix de la methodologie de travail.

Le \essentiel{deuxieme chapitre} conduit l'analyse globale. Il identifie
les acteurs du systeme, recense les besoins fonctionnels et non
fonctionnels, les modelise par un diagramme de cas d'utilisation global,
puis presente le pilotage du projet : composition de l'equipe, carnet de
produit et decoupage en sprints.

Le \essentiel{troisieme chapitre} decrit l'environnement de travail,
materiel et logiciel, puis l'architecture retenue. Il en donne les deux
vues complementaires, logique et physique, et explique les trois choix
structurants qui la caracterisent.

Les \essentiel{chapitres quatre a neuf} rendent compte des six sprints.
Chacun suit la meme trame : planification et carnet de sprint,
elaboration technique appuyee sur les modeles UML, realisation et tests
de validation.

La \essentiel{conclusion generale} dresse le bilan du travail accompli,
revient sur les difficultes rencontrees, evalue l'atteinte des objectifs
et ouvre sur les perspectives d'evolution de la plateforme.


% =====================================================================
%  Chapitre 1 - Etude prealable
% =====================================================================
\chapter{Etude prealable}
\label{chap:etude-prealable}

\section*{Introduction}
\addcontentsline{toc}{section}{Introduction}

Tout projet informatique repond a un besoin ne d'une situation concrete.
Ce premier chapitre etablit cette situation. Il presente l'organisme au
sein duquel le projet a ete conduit, formule la problematique a laquelle
il entend repondre, examine la maniere dont ce besoin est aujourd'hui
couvert et en releve les limites. Il expose ensuite la solution retenue
et justifie la methodologie de travail adoptee.

% =====================================================================
\section{Presentation de l'organisme d'accueil}

\subsection{L'entreprise}

Le present projet de fin d'etudes a ete realise au sein de
\essentiel{\entreprise}, entreprise de services numeriques implantee a
\entrepriseVille{} (\entrepriseCodePostal), en Tunisie.

L'activite de \entreprise{} s'organise autour de trois metiers
complementaires.

\begin{table}[H]
\centering
\caption{Les trois metiers de \entreprise}
\label{tab:metiers-sipro}
\small
\begin{tabularx}{\textwidth}{@{}p{3.6cm}X@{}}
\toprule
\textbf{Metier} & \textbf{Perimetre} \\ \midrule
Edition logicielle & Logiciels de gestion et progiciels d'entreprise.
  Metier historique, dont releve le present projet \\
Integration RFID & Identification par radiofrequence : tracabilite des
  biens, inventaire, controle d'acces \\
Materiel informatique & Vente, maintenance et assistance sur le materiel
  professionnel \\
\bottomrule
\end{tabularx}
\end{table}

Cette combinaison eclaire le projet. Une entreprise qui integre des
solutions de tracabilite et maintient du materiel sur site connait de
l'interieur les contraintes de la gestion d'un parc : suivi des
equipements, planification des interventions, surveillance des echeances.
\essentiel{Le sujet confie prolonge donc un savoir-faire existant vers un
domaine ou la demande des clients de \entreprise{} etait exprimee.}

\subsection{La mission confiee}

\entreprise{} a confie la conception et le developpement de
\essentiel{\sujetCourt}, une plateforme web complete de gestion de flotte
et de transport de marchandises, destinee a digitaliser les operations
des societes de transport clientes de l'entreprise.

Le stage s'est deroule du \stageDebut{} au \stageFin, soit une duree de
\stageDuree{} mois. Cette duree a rendu possible un decoupage en six
iterations d'un mois environ, chacune consacree a un domaine metier
coherent et livrant un increment utilisable.

Les objectifs fixes au depart etaient au nombre de quatre.

\begin{enumerate}
  \item \essentiel{Centraliser} dans un referentiel unique les vehicules,
    les equipements, les personnels et les clients, aujourd'hui disperses
    entre plusieurs documents.
  \item \essentiel{Industrialiser} la maintenance : ordres de travail,
    interventions, stock de pieces de rechange et surveillance des
    echeances reglementaires.
  \item \essentiel{Ouvrir un canal numerique au client final}, de la
    demande de transport au suivi de l'acheminement et a la consultation
    des factures.
  \item \essentiel{Outiller la decision} par l'estimation tarifaire,
    l'anticipation des pannes et l'affectation assistee des ressources.
\end{enumerate}

Le livrable attendu etait une plateforme fonctionnelle, deployable et
documentee, accompagnee de la modelisation UML correspondante.

% =====================================================================
\section{Problematique}
\label{sec:problematique}

Une entreprise de transport de marchandises pilote simultanement quatre
ressources dont les cycles de vie sont independants : un parc de
vehicules, un ensemble d'equipements embarques, une equipe de chauffeurs
et un portefeuille de clients. Chacune produit ses propres evenements,
et c'est de leur croisement que naissent les difficultes.

\figReelle{problematique.png}
  {Interdependance des ressources d'une entreprise de transport}
  {fig:problematique}{0.55}

La figure~\ref{fig:problematique} illustre cette interdependance. Elle se
traduit concretement par cinq difficultes recurrentes.

\paragraph{La dispersion de l'information.}
Le suivi des vehicules, celui des interventions de maintenance, le
planning des chauffeurs et la facturation vivent dans des documents
distincts, souvent detenus par des personnes differentes. Repondre a une
question aussi simple que \exemple{ce vehicule est-il disponible jeudi
prochain ?} suppose de consulter plusieurs sources et d'esperer qu'elles
soient a jour.

\paragraph{L'absence d'anticipation sur la maintenance.}
Les interventions sont declenchees par la panne plutot que par
l'observation de son approche. Le cout d'une immobilisation imprevue
depasse largement celui de l'entretien qui l'aurait evitee : il faut y
ajouter la mission non honoree, le report sur un autre vehicule et
l'insatisfaction du client.

\paragraph{Le risque administratif.}
Une visite technique expiree, une assurance echue ou un permis de
conduire perime exposent l'entreprise a une immobilisation immediate et a
des sanctions. Ces echeances sont nombreuses, dispersees dans le temps et
rarement surveillees de maniere systematique.

\paragraph{L'estimation tarifaire empirique.}
Le prix d'un transport se fixe souvent d'apres l'experience de la
personne qui etablit le devis. Deux devis comparables peuvent ainsi
differer sensiblement, ce qui fragilise la coherence commerciale et rend
la marge difficile a piloter.

\paragraph{L'opacite pour le client.}
Le client n'a pas de moyen autonome de deposer une demande, de suivre
l'acheminement de sa marchandise ni de consulter ses factures. Chaque
demande d'information mobilise un interlocuteur, ce qui alourdit le
travail administratif des deux cotes.

\vspace{0.4cm}
\noindent
\essentiel{Ces cinq difficultes ont une cause commune : l'absence d'un
referentiel unique ou les vehicules, les equipements, les chauffeurs, les
missions et les clients seraient decrits une seule fois et relies entre
eux.} C'est ce constat qui fonde le projet.

% =====================================================================
\section{Etude de l'existant}
\label{sec:existant}

\subsection{La situation avant le projet}

Chez les societes de transport clientes de \entreprise, la gestion
reposait sur des \essentiel{methodes traditionnelles}, articulees autour
de trois supports.

\begin{itemize}
  \item \essentiel{Des feuilles de calcul Excel} pour le suivi du parc,
    des missions et de la facturation. Chaque classeur etait tenu par une
    personne, sur son poste, dans sa propre organisation de colonnes.
  \item \essentiel{Des echanges par appel telephonique et par WhatsApp}
    pour la coordination quotidienne : affectation d'un chauffeur,
    signalement d'une panne, confirmation d'une livraison. Ces echanges
    portaient l'information operationnelle mais n'en laissaient aucune
    trace exploitable.
  \item \essentiel{Un tableau de bord rudimentaire}, alimente
    manuellement, offrant une vue partielle et rarement a jour.
\end{itemize}

Cette organisation fonctionne tant que le volume reste modeste et que les
memes personnes restent en poste. Elle atteint sa limite des que
plusieurs interlocuteurs doivent consulter ou modifier la meme
information, et surtout des qu'il faut \exemple{retrouver ce qui a ete
decide il y a trois mois}.

\subsection{Les alternatives du marche}

Face a cette situation, trois familles d'outils sont envisageables.

\paragraph{Le tableur.}
C'est la situation de depart decrite ci-dessus. Le tableur a pour lui sa
gratuite, sa souplesse et le fait que chacun sache s'en servir. Il ne
gere en revanche ni les acces concurrents, ni la tracabilite des
modifications, ni les regles metier.

\paragraph{Le progiciel de gestion generaliste.}
Certaines entreprises s'appuient sur un progiciel comptable ou commercial
auquel elles greffent le suivi de la flotte. La comptabilite et la
facturation y sont bien traitees, mais les notions propres au transport
(mission, echeance reglementaire, ordre de travail, piece de rechange)
n'existent pas nativement et doivent etre detournees d'autres objets.

\paragraph{La solution specialisee en ligne.}
Il existe des offres dediees a la gestion de flotte, generalement
proposees en abonnement. Elles couvrent bien le suivi des vehicules et la
geolocalisation, mais leur cout est proportionnel a la taille du parc,
leur parametrage suppose une adaptation aux usages du fournisseur, et
l'hebergement des donnees echappe a l'entreprise.

\subsection{Critique de l'existant}

Le tableau~\ref{tab:critique-existant} confronte ces trois familles aux
besoins identifies dans la section~\ref{sec:problematique}.

\begin{table}[H]
\centering
\caption{Couverture des besoins par les pratiques existantes}
\label{tab:critique-existant}
\small
\begin{tabularx}{\textwidth}{@{}Xccc@{}}
\toprule
\textbf{Besoin} & \textbf{Tableur} & \textbf{Progiciel} & \textbf{Solution} \\
                &                  & \textbf{generaliste} & \textbf{specialisee} \\
\midrule
Referentiel unique des ressources      & non      & partiel  & oui \\
Travail simultane a plusieurs          & non      & oui      & oui \\
Tracabilite des modifications          & non      & oui      & oui \\
Surveillance des echeances reglementaires & non   & non      & partiel \\
Maintenance preventive et ordres de travail & non & non      & partiel \\
Estimation tarifaire outillee          & non      & non      & non \\
Portail autonome pour le client        & non      & non      & partiel \\
Suivi de mission en direct             & non      & non      & oui \\
Maitrise des donnees et du cout        & oui      & partiel  & non \\
\bottomrule
\end{tabularx}
\end{table}

Aucune des trois familles ne couvre l'ensemble du besoin. Les deux
manques les plus nets concernent \essentiel{l'aide a la decision}
(estimation tarifaire, anticipation des pannes, affectation optimisee des
ressources), absente partout, et \essentiel{la surveillance automatique
des echeances}, qui n'est traitee qu'en partie et rarement de facon
transversale aux vehicules, aux chauffeurs et aux equipements.

\subsection{Solution proposee}

\sujetCourt{} est une plateforme unifiee de gestion de flotte et de
transport de marchandises. Elle repose sur trois principes.

\paragraph{Un referentiel unique, des domaines separes.}
Chaque ressource est decrite une seule fois. Les domaines metier restent
toutefois autonomes : la facturation, la maintenance et l'exploitation
evoluent independamment, ce qui permet de faire progresser un domaine
sans immobiliser les autres.

\paragraph{Trois publics, trois applications.}
Un back-office pour le personnel de l'entreprise, un portail pour les
clients, une application mobile pour les chauffeurs. Chacune n'expose que
ce qui concerne son public.

\paragraph{L'aide a la decision integree au flux de travail.}
L'estimation tarifaire, l'anticipation des pannes, l'affectation des
ressources et la veille metier ne sont pas des modules a part : elles
interviennent au moment ou la decision se prend, dans l'ecran ou
l'utilisateur travaille.

\begin{definitionbox}{Ce que couvre la plateforme}
Douze services metier autonomes, trois applications destinees a des
publics distincts, un referentiel commun de vehicules, d'equipements,
de personnels et de clients, et quatre services d'aide a la decision.
Le detail de ce perimetre est etabli au chapitre suivant.
\end{definitionbox}

% =====================================================================
\section{Methodologie de travail}
\label{sec:methodologie}

Le choix d'une methodologie n'est pas neutre : il determine la maniere
dont les besoins sont recueillis, dont le travail est decoupe et dont
l'avancement est mesure.

\subsection{Le developpement en cascade}

Le modele en cascade enchaine des phases sequentielles : specification,
conception, developpement, tests, mise en service. Chaque phase se
termine avant que la suivante ne commence, et produit un document qui
sert de contrat pour la suite.

Ce modele convient aux projets dont le besoin est stable et connu des le
depart. Il presente en revanche deux inconvenients pour un projet comme
celui-ci : le client ne voit un resultat utilisable qu'a la toute fin, et
une erreur de comprehension du besoin ne se decouvre qu'au moment des
tests, lorsque son cout de correction est maximal.

\subsection{L'approche agile}

Les approches agiles renversent cette logique. Le travail est decoupe en
iterations courtes, chacune produisant un increment fonctionnel du
produit. Le besoin est precise au fil de l'eau, au contact du resultat
plutot que sur document.

\subsection{Le cadre Scrum}

Scrum est le cadre agile le plus repandu. Il repose sur trois roles, sur
des evenements qui rythment l'iteration et sur trois artefacts.

\figReelle{methodo-scrum.jpg}
  {Le processus Scrum}
  {fig:scrum}{0.72}

\begin{table}[H]
\centering
\caption{Roles, evenements et artefacts de Scrum}
\label{tab:scrum}
\small
\begin{tabularx}{\textwidth}{@{}p{2.8cm}X@{}}
\toprule
\textbf{Element} & \textbf{Definition} \\ \midrule
\textit{Roles} & Product Owner : porte la valeur metier et ordonne le
  carnet de produit. Scrum Master : garantit le cadre et leve les
  obstacles. Equipe de developpement : concoit, developpe et teste
  l'increment \\
\textit{Evenements} & Planification, melee quotidienne, revue de sprint,
  retrospective \\
\textit{Artefacts} & Carnet de produit, carnet de sprint, increment \\
\bottomrule
\end{tabularx}
\end{table}

\subsection{Choix retenu et adaptation au contexte}

\essentiel{Scrum a ete retenu pour ce projet}, pour trois raisons.

D'abord, \essentiel{le besoin n'etait pas entierement connu au depart}.
Le perimetre s'est precise au contact des utilisateurs : plusieurs
fonctionnalites, comme l'archivage reversible des comptes ou la
surveillance des echeances documentaires, ont ete identifiees en cours de
projet et n'auraient pas figure dans une specification initiale.

Ensuite, \essentiel{le decoupage en domaines metier se prete
naturellement a l'iteration}. Chaque sprint peut livrer un domaine
complet et utilisable, ce qui rend l'avancement lisible et permet un
retour utilisateur rapide.

Enfin, \essentiel{le format du projet impose un rythme}. Un stage de fin
d'etudes se deroule sur une duree fixe : le decoupage en iterations de
duree egale donne une mesure objective de l'avancement et permet de
reajuster le perimetre en cours de route plutot qu'a la fin.

\paragraph{Adaptation au contexte du stage.}
Scrum suppose une equipe de plusieurs developpeurs. Le projet ayant ete
mene par un seul developpeur, le cadre a ete adapte : les roles de
Product Owner et de Scrum Master ont ete tenus par les encadrants, la
melee quotidienne est devenue un point d'avancement hebdomadaire, et la
revue de sprint a pris la forme d'une demonstration a l'encadrant
entreprise. Les artefacts, en revanche, ont ete tenus integralement :
carnet de produit, carnet de sprint et increment livre a chaque
iteration.

\subsection{Outils de pilotage}
\label{sec:outils-pilotage}

Trois outils ont soutenu la conduite du projet au quotidien.

\begin{table}[H]
\centering
\caption{Outils de pilotage du projet}
\label{tab:outils-pilotage}
\small
\begin{tabularx}{\textwidth}{@{}p{2.2cm}p{3.4cm}X@{}}
\toprule
\textbf{Outil} & \textbf{Usage} & \textbf{Apport} \\ \midrule
Trello & Carnet et suivi des taches & Chaque recit devient une carte,
  deplacee au fil de son avancement : le carnet de sprint est visible
  d'un coup d'oeil \\
Git et GitHub & Gestion de versions & Historique, branches, retour
  arriere, sauvegarde et support aux revues de code \\
WhatsApp & Communication quotidienne & Canal court avec l'encadrant
  entreprise, sans attendre le point hebdomadaire \\
\bottomrule
\end{tabularx}
\end{table}

\begin{figure}[H]
  \centering
  \includegraphics[height=1.35cm,keepaspectratio]{trello-logo.jpg}\hspace{1.4cm}
  \includegraphics[height=1.35cm,keepaspectratio]{logo-git.png}\hspace{1.4cm}
  \includegraphics[height=1.35cm,keepaspectratio]{whatsapplogo.jpg}
  \caption{Outils de pilotage et de communication du projet}
  \label{fig:outils-pilotage}
\end{figure}

\essentiel{WhatsApp figure a la fois dans l'existant critique et dans les
outils de pilotage.} La contradiction n'est qu'apparente : ce qui pose
probleme n'est pas l'outil mais son emploi comme \exemple{support
d'information metier}. Coordonner une equipe par messagerie instantanee
est legitime ; y conserver l'etat d'un parc de vehicules ne l'est pas.
C'est la frontiere que la plateforme etablit.

% =====================================================================
\section*{Conclusion}
\addcontentsline{toc}{section}{Conclusion}

Ce chapitre a etabli le contexte du projet. L'examen des pratiques
existantes a montre qu'aucune des trois familles d'outils rencontrees ne
couvre l'ensemble du besoin d'une entreprise de transport de
marchandises, les manques les plus nets portant sur l'aide a la decision
et sur la surveillance des echeances reglementaires. La solution
proposee, \sujetCourt, repose sur un referentiel unique servant des
domaines metier autonomes, sur trois applications adaptees a trois
publics et sur une aide a la decision integree au flux de travail. Le
cadre Scrum a ete retenu pour conduire cette realisation, adapte au
contexte d'un stage, et outille par Trello, Git et GitHub.

Le chapitre suivant traduit ce cadrage en exigences : il identifie les
acteurs du systeme, recense les besoins fonctionnels et non fonctionnels,
et etablit le decoupage du travail en six sprints.


% Les chapitres suivants sont rediges un par un, apres validation du
% precedent. Decommenter au fur et a mesure.
% =====================================================================
%  Chapitre 2 - Etude et analyse globale du projet
% =====================================================================
\chapter{Etude et analyse globale du projet}
\label{chap:analyse}

\section*{Introduction}
\addcontentsline{toc}{section}{Introduction}

Le chapitre precedent a etabli le besoin. Celui-ci le traduit en
exigences. Il identifie d'abord les acteurs qui interagissent avec la
plateforme, puis recense les besoins fonctionnels domaine par domaine et
les besoins non fonctionnels qui conditionnent la qualite du systeme. Ces
exigences sont ensuite synthetisees par un diagramme de cas d'utilisation
global. Le chapitre se clot sur le pilotage du projet : composition de
l'equipe, carnet de produit et decoupage du travail en six sprints.

% =====================================================================
\section{Identification des acteurs}
\label{sec:acteurs}

Un acteur est une entite exterieure au systeme qui interagit avec lui.
L'analyse des habilitations reellement portees par la plateforme fait
apparaitre six acteurs humains et deux systemes tiers.

\subsection{Les acteurs principaux}

\begin{table}[H]
\centering
\caption{Acteurs principaux de la plateforme}
\label{tab:acteurs}
\small
\begin{tabularx}{\textwidth}{@{}p{2.4cm}X@{}}
\toprule
\textbf{Acteur} & \textbf{Role et perimetre d'action} \\ \midrule
Utilisateur & Acteur generique dont derivent tous les autres : compte,
  profil, second facteur, consultations non restreintes, messagerie \\
Client & Donneur d'ordre, entreprise ou particulier. Depose une demande,
  recoit un devis, suit l'acheminement, consulte ses factures, ouvre des
  tickets de support \\
Employe & Acteur \essentiel{abstrait} regroupant le personnel interne.
  Il n'agit pas : il factorise ce que partagent les roles internes \\
Chauffeur & Conduit les vehicules. Connexion par code PIN, consultation
  de ses missions, transmission de sa position \\
Technicien & Intervient sur les vehicules : interventions, consommation
  de pieces, avancement des ordres de travail, documentation \\
Responsable & Acteur \essentiel{abstrait} couvrant la combinaison
  Administrateur ou Gestionnaire, la plus frequente du code. Referentiels,
  devis, missions, maintenance, facturation \\
Administrateur & Droits exclusifs : approbation et rejet des comptes,
  habilitations, archivage, purge, suppressions definitives \\
Gestionnaire & Exploitation courante, sans droit de suppression ni
  d'administration des comptes \\
\bottomrule
\end{tabularx}
\end{table}

\subsection{La hierarchie des acteurs}

Ces acteurs ne sont pas independants : ils forment une hierarchie de
generalisation calquee sur l'heritage reel du modele de donnees, ou la
table des utilisateurs porte un discriminant distinguant les clients des
employes.

\begin{definitionbox}{Pourquoi une hierarchie d'acteurs}
Sans generalisation, un cas d'utilisation ouvert a tous devrait etre
relie separement au Client, au Chauffeur, au Technicien, au Gestionnaire
et a l'Administrateur, soit cinq associations pour une seule capacite.
La generalisation permet de ne le relier qu'a l'acteur \exemple{le plus
general qui le possede} : les acteurs specialises en heritent. Le
diagramme gagne en lisibilite sans rien perdre en precision.
\end{definitionbox}

Deux acteurs de cette hierarchie sont \essentiel{abstraits} : Employe et
Responsable ne correspondent a aucun role stocke en base. Ils
n'existent que pour factoriser. \essentiel{Responsable} merite une
explication : il traduit la combinaison d'habilitations
\exemple{Administrateur ou Gestionnaire}, qui revient dans la majorite
des controles d'acces de la plateforme. Le nommer evite de dupliquer
chaque association.

% =====================================================================
\section{Besoins fonctionnels}
\label{sec:besoins-fonctionnels}

Les besoins fonctionnels sont presentes par domaine metier. Ce decoupage
n'est pas seulement un procede de redaction : c'est celui qui structure
la plateforme et, on le verra en section~\ref{sec:sprints}, celui qui
determine le decoupage en sprints.

\begin{table}[H]
\centering
\caption{Besoins fonctionnels par domaine metier}
\label{tab:besoins-fonctionnels}
\small
\begin{tabularx}{\textwidth}{@{}p{3.1cm}X@{}}
\toprule
\textbf{Domaine} & \textbf{Besoins recenses} \\ \midrule

Identites, acces et securite &
  S'inscrire et confirmer son adresse electronique ; se connecter, avec
  validation d'un second facteur lorsqu'il est actif ; se connecter par
  code PIN depuis l'application chauffeur ; gerer son profil et son mot
  de passe ; activer ou desactiver son second facteur ; statuer sur les
  demandes d'inscription ; gerer les habilitations ; archiver puis
  restaurer un compte ; consulter et effacer l'historique des
  connexions \\

Referentiel de flotte et d'equipements &
  Consulter le parc et rechercher un vehicule disponible ; tenir la fiche
  d'un vehicule ; gerer l'inventaire des equipements, leur affectation et
  leur etat ; enregistrer les pieces administratives et les assurances et
  surveiller leurs echeances ; affecter et restituer un equipement ;
  realiser une inspection ; declarer les pleins et suivre la
  consommation ; sortir un bien de l'inventaire \\

Demande, devis, mission et suivi &
  Deposer une demande depuis le portail client avec estimation
  tarifaire ; etablir un devis en retenant, modifiant ou fixant le prix ;
  accepter ou refuser un devis, la mission n'etant creee qu'a
  l'acceptation ; affecter un vehicule et un chauffeur, au besoin avec
  l'assistance du moteur d'optimisation ; transmettre la position depuis
  l'application mobile ; suivre une mission en direct et consulter
  l'historique du trajet ; cloturer ou annuler une mission ; ouvrir et
  traiter un ticket de support \\

Maintenance curative et predictive &
  Ouvrir, instruire et valider un ordre de travail et en suivre
  l'avancement par etat ; declarer une intervention, la suivre, consommer
  ou restituer des pieces ; tenir le catalogue des pieces et
  reapprovisionner le stock avec alerte sur les seuils bas ; documenter
  un ordre de travail ; consulter le risque de panne par vehicule et
  l'organe concerne \\

Facturation, alertes et journal metier &
  Etablir une facture a partir des missions facturables, avec controle
  des doublons ; suivre le cycle de vie d'une facture et identifier les
  impayes ; consulter ses factures depuis le portail client ; recevoir
  les alertes metier portant sur les echeances, les stocks, les retards
  et les impayes ; consulter le journal des evenements et editer un
  rapport de synthese \\

Collaboration et aide a la decision &
  Echanger par messagerie instantanee, en conversation privee ou en
  groupe, avec pieces jointes ; consulter les correspondants connectes ;
  interroger un assistant documentaire sur les services et les donnees ;
  consulter un tableau de bord consolidant les indicateurs
  d'exploitation \\

\bottomrule
\end{tabularx}
\end{table}

Ces six domaines ne sont pas un simple classement de redaction :
\essentiel{ils determinent le decoupage en sprints} presente en
section~\ref{sec:sprints}, et chacun donnera son chapitre a la partie
consacree a la realisation.

% =====================================================================
\section{Besoins non fonctionnels}
\label{sec:besoins-non-fonctionnels}

Les besoins non fonctionnels ne decrivent pas ce que fait le systeme mais
la maniere dont il doit le faire. Ils ont oriente les choix
d'architecture presentes au chapitre suivant.

\begin{table}[H]
\centering
\caption{Besoins non fonctionnels et reponses apportees}
\label{tab:besoins-nf}
\small
\begin{tabularx}{\textwidth}{@{}p{2.6cm}X@{}}
\toprule
\textbf{Exigence} & \textbf{Contenu et reponse retenue} \\ \midrule
Securite & Jeton signe, second facteur optionnel, controle d'acces par
  role, verrouillage apres echecs repetes, journalisation des connexions.
  Aucune suppression definitive d'un compte porteur d'historique \\
Tracabilite & Toute operation sensible laisse une trace : historique des
  statuts d'ordre de travail, journal des connexions, journal metier.
  Reponse directe a la volatilite des echanges informels \\
Performance & Pagination systematique et filtrage cote serveur : les
  ecrans restent utilisables sur plusieurs milliers de lignes \\
Disponibilite & Chaque service expose un point de controle de sante ;
  un service defaillant n'empeche pas les autres de fonctionner \\
Maintenabilite & Responsabilites separees dans chaque service, services
  independants entre eux : faire evoluer la facturation n'oblige pas a
  redeployer la maintenance \\
Portabilite & Deploiement par conteneurs, sans installation manuelle. Le
  poste et le serveur executent la meme definition \\
Ergonomie & Trois interfaces plutot qu'une : chaque public ne voit que ce
  qui le concerne. Le chauffeur dispose d'une application mobile dediee \\
Evolutivite & Ajouter un domaine ne doit pas imposer de reprendre
  l'existant. Le huitieme service, la maintenance predictive, a ete
  ajoute apres coup : il valide ce point \\
\bottomrule
\end{tabularx}
\end{table}

% =====================================================================
\section{Modelisation des besoins : cas d'utilisation global}
\label{sec:uc-global}

Le diagramme de la figure~\ref{fig:uc-global} synthetise les besoins
fonctionnels. Il donne la carte d'ensemble de la plateforme ; le detail
de chaque domaine est presente dans le chapitre du sprint correspondant.

\figReelle{platform-usecase-global.png}
  {Diagramme de cas d'utilisation global de la plateforme \sujetCourt}
  {fig:uc-global}{0.60}

\paragraph{Lecture du diagramme.}
Les acteurs sont alignes a gauche, relies entre eux par des liens de
generalisation dont le triangle creux pointe vers l'acteur le plus
general. La frontiere du systeme est subdivisee en six domaines
fonctionnels, qui correspondent aux six sections des besoins
fonctionnels. Les deux systemes tiers figurent a droite, hors de la
frontiere.

\paragraph{Les cas d'aide a la decision.}
Trois cas sont volontairement places hors des domaines :
\exemple{estimer le prix du trajet}, \exemple{optimiser l'affectation de
la flotte} et \exemple{anticiper les pannes}. Ils ne constituent pas un
but d'utilisateur en eux-memes : personne ne se connecte pour
\exemple{lancer une prediction}. Ils sont atteints \essentiel{depuis un
cas de base}, au moment ou la decision se prend. C'est ce que traduisent
les relations d'inclusion et d'extension :

\begin{itemize}
  \item deposer une demande de transport \essentiel{inclut}
    systematiquement l'estimation du prix : la fourchette tarifaire est
    affichee au client avant qu'il ne valide ;
  \item etablir un devis est \essentiel{etendu} par l'estimation :
    le gestionnaire peut retenir le prix propose, le modifier ou le fixer
    lui-meme, l'estimation reste facultative ;
  \item planifier une mission \essentiel{inclut} l'optimisation, et
    piloter les ordres de travail \essentiel{inclut} l'anticipation des
    pannes.
\end{itemize}

\paragraph{Ce que le diagramme ne montre pas.}
Les services d'aide a la decision sont ici \essentiel{internes} au
systeme : ce sont des composants de la plateforme, pas des entites
exterieures. Seuls le serveur de messagerie et la passerelle de
messagerie instantanee, qui appartiennent a des tiers, figurent comme
acteurs secondaires.

% =====================================================================
\section{Pilotage du projet avec Scrum}
\label{sec:pilotage}

\subsection{Equipe et roles}

Le cadre Scrum a ete transpose au contexte d'un projet de fin d'etudes
selon la repartition du tableau~\ref{tab:equipe}.

\begin{table}[H]
\centering
\caption{Equipe projet et repartition des roles Scrum}
\label{tab:equipe}
\small
\begin{tabularx}{\textwidth}{@{}p{3.4cm}p{4.2cm}X@{}}
\toprule
\textbf{Role Scrum} & \textbf{Tenu par} & \textbf{Responsabilites} \\ \midrule
Product Owner & \encadrantEntreprise{} \newline (encadrant entreprise) &
  Porte le besoin metier, ordonne le carnet, valide les increments \\
Scrum Master & \encadrantAcademique{} \newline (encadrant academique) &
  Garantit le cadre et la rigueur methodologique, leve les obstacles \\
Equipe de developpement & \etudiantComplet &
  Conception, developpement, tests et deploiement \\
\bottomrule
\end{tabularx}
\end{table}

\subsection{Carnet de produit}

Le carnet de produit rassemble les besoins exprimes sous forme de recits
utilisateur, ordonnes par priorite. La priorisation suit la methode
MoSCoW : \exemple{Must have} pour ce qui est indispensable,
\exemple{Should have} pour ce qui est important sans etre bloquant,
\exemple{Could have} pour ce qui apporte un confort. L'estimation est
exprimee en points, selon une suite de Fibonacci.

\begin{table}[H]
\centering
\caption{Carnet de produit, par theme}
\label{tab:backlog-produit}
\small
\begin{tabularx}{\textwidth}{@{}p{0.9cm}Xp{1.9cm}p{1.4cm}@{}}
\toprule
\textbf{Ref.} & \textbf{Theme} & \textbf{Priorite} & \textbf{Points} \\
\midrule
T1 & Identites, acces et securite            & Must have   & 34 \\
T2 & Referentiel de flotte et d'equipements  & Must have   & 34 \\
T3 & Demande, devis, mission et suivi        & Must have   & 55 \\
T4 & Maintenance curative                    & Must have   & 34 \\
T5 & Facturation                             & Must have   & 21 \\
T6 & Alertes metier                          & Should have & 13 \\
T7 & Messagerie interne                      & Should have & 21 \\
T8 & Estimation tarifaire assistee           & Should have & 13 \\
T9 & Suivi de position en direct             & Should have & 21 \\
T10 & Maintenance predictive                 & Could have  & 13 \\
T11 & Journal metier et rapports             & Could have  & 13 \\
T12 & Assistant documentaire                 & Could have  & 21 \\
T13 & Optimisation de l'affectation          & Could have  & 13 \\
\midrule
\multicolumn{3}{@{}r}{\textbf{Total}} & \textbf{306} \\
\bottomrule
\end{tabularx}
\end{table}

\subsection{Decoupage en sprints}
\label{sec:sprints}

Le stage s'etendant sur \stageDuree{} mois, le travail a ete decoupe en
six sprints, precedes d'une phase de cadrage et suivis d'une phase de
recette. Chaque sprint couvre un \essentiel{domaine metier coherent} et
livre un increment utilisable de bout en bout : interface, service,
persistance.

\begin{table}[H]
\centering
\caption{Contenu des six sprints}
\label{tab:sprints}
\small
% La colonne « Duree » valait 1,2 cm : « 4 semaines » n'y tenait pas et
% se coupait en « 4 se-maines ». Portee a 2,3 cm, prise sur la colonne
% flexible qui a de la marge.
\begin{tabularx}{\textwidth}{@{}p{1.5cm}Xp{2.6cm}p{2.3cm}@{}}
\toprule
\textbf{Sprint} & \textbf{Domaine couvert} & \textbf{Themes} & \textbf{Duree} \\
\midrule
Sprint 1 & Identites, acces et securite & T1 & 4 semaines \\
Sprint 2 & Referentiel de flotte et d'equipements & T2 & 4 semaines \\
Sprint 3 & Demande, devis, mission, suivi en direct & T3, T8, T9 & 4 semaines \\
Sprint 4 & Maintenance curative et predictive & T4, T10 & 4 semaines \\
Sprint 5 & Facturation, alertes et journal metier & T5, T6, T11 & 3 semaines \\
Sprint 6 & Collaboration et aide a la decision & T7, T12, T13 & 3 semaines \\
\bottomrule
\end{tabularx}
\end{table}

L'ordre n'est pas arbitraire. \essentiel{Le sprint 1 conditionne tous les
autres} : sans identites ni habilitations, aucun ecran ne peut etre
protege. \essentiel{Le sprint 2 fournit le referentiel} sur lequel
s'appuient les missions et la maintenance : on ne peut affecter un
vehicule qui n'existe pas. Les sprints suivants exploitent ces
fondations, ce qui explique que les themes d'aide a la decision soient
traites au sein du sprint du domaine qu'ils servent, et non isolement.

\figReelle{planning-sprints.png}
  {Planification previsionnelle du projet}
  {fig:planning}{1.00}

Le diagramme de Gantt de la figure~\ref{fig:planning} situe ces phases
dans le calendrier du stage, du \stageDebut{} au \stageFin. La redaction
du rapport y apparait comme une activite transversale, engagee des le
quatrieme sprint plutot que reportee en fin de projet.

% =====================================================================
\section*{Conclusion}
\addcontentsline{toc}{section}{Conclusion}

Ce chapitre a traduit le besoin en exigences. Six acteurs humains
organises en hierarchie de generalisation et deux systemes tiers
interagissent avec la plateforme. Les besoins fonctionnels, presentes par
domaine, ont ete synthetises par un diagramme de cas d'utilisation global
dont la lecture met en evidence le parti pris principal : l'aide a la
decision n'est pas un module separe mais un ensemble de cas atteints
depuis les cas metier. Les besoins non fonctionnels ont ete formules et
rattaches aux reponses qui leur seront apportees. Le pilotage, enfin, a
ete arrete : trois roles Scrum, un carnet de produit de treize themes et
un decoupage en six sprints dont l'ordre suit les dependances entre
domaines.

Le chapitre suivant presente l'environnement de travail et l'architecture
qui repond a ces exigences.

% =====================================================================
%  Chapitre 3 - Environnement de travail et architecture
% =====================================================================
\chapter{Environnement de travail et architecture}
\label{chap:architecture}

\section*{Introduction}
\addcontentsline{toc}{section}{Introduction}

Les exigences etant posees, ce chapitre presente les moyens mis en oeuvre
pour y repondre. Il decrit d'abord l'environnement de travail, materiel
et logiciel, puis l'architecture de la plateforme sous ses deux vues
complementaires : la vue logique, qui organise les responsabilites, et la
vue physique, qui decrit le deploiement. Il se termine par l'examen des
trois choix structurants qui caracterisent cette architecture, y compris
ceux qui s'ecartent du patron de reference.

% =====================================================================
\section{Environnement materiel}

Le developpement a ete conduit sur un poste unique, qui heberge
egalement l'ensemble des conteneurs de la plateforme.

\begin{table}[H]
\centering
\caption{Caracteristiques du poste de developpement}
\label{tab:materiel}
\small
\begin{tabularx}{\textwidth}{@{}p{3.4cm}X@{}}
\toprule
\textbf{Element} & \textbf{Caracteristique} \\ \midrule
Systeme d'exploitation & Windows 11 Professionnel \\
Environnement conteneurise & Docker Desktop, avec le sous-systeme Linux \\
Memoire vive & 16 Go, dont environ 8 Go alloues aux conteneurs \\
Stockage & Disque a semi-conducteurs, indispensable aux temps de
  compilation des huit services \\
\bottomrule
\end{tabularx}
\end{table}

Ce point merite d'etre souligne : \essentiel{une architecture
microservices se developpe difficilement sans conteneurisation}. Vingt
conteneurs cohabitent, dont une base de donnees, un courtier de messages
et une base vectorielle. Les installer et les configurer un a un sur le
poste serait long, fragile et non reproductible.

% =====================================================================
\section{Environnement logiciel}

\subsection{Outils de developpement et de gestion de version}

\begin{table}[H]
\centering
\caption{Outils utilises}
\label{tab:outils}
\small
\begin{tabularx}{\textwidth}{@{}p{3.2cm}p{3.0cm}X@{}}
\toprule
\textbf{Outil} & \textbf{Usage} & \textbf{Motif du choix} \\ \midrule
Visual Studio Code & Edition du code &
  Un seul editeur pour C\#, TypeScript, Python et SQL. \\
Docker Desktop & Execution locale &
  La definition des services est versionnee : le poste et le serveur
  executent la meme chose. \\
Git et GitHub & Gestion de version &
  Historique, branches et sauvegarde. \\
Postman & Essais d'API &
  Verification des points d'entree REST avant toute integration
  d'interface. \\
pgAdmin & Administration de la base &
  Inspection du schema et des donnees. \\
PlantUML & Modelisation UML &
  Les diagrammes sont du texte : ils se versionnent et se comparent
  comme du code. \\
Trello & Suivi du carnet &
  Voir la section~\ref{sec:outils-pilotage}. \\
\bottomrule
\end{tabularx}
\end{table}

\subsection{Langages, cadriciels et briques techniques}

\begin{table}[H]
\centering
\caption{Technologies retenues et versions}
\label{tab:technologies}
\small
\begin{tabularx}{\textwidth}{@{}p{2.9cm}p{2.6cm}X@{}}
\toprule
\textbf{Domaine} & \textbf{Briques et versions} & \textbf{Role} \\
\midrule
Services metier & .NET / C\# 8.0, ASP.NET Core 8.0, Entity Framework
  Core et Npgsql 8.0.11, MassTransit 8.3.6, SignalR, BCrypt.Net 4.0.3,
  FluentValidation 11.9.0, Serilog 8.0.3 &
  Socle des huit services metier, exposition REST, correspondance
  objet-relationnel, messagerie et journalisation \\
Interfaces & React 18.3 et Vite 6.0, Next.js 14.2, Expo 54 avec React
  Native 0.81, TypeScript 5.6, Leaflet 1.9, ApexCharts 4.4 &
  Back-office, portail client, application chauffeur, cartographie et
  graphiques \\
Aide a la decision & Python 3.11 et 3.12, FastAPI 0.115, XGBoost 2.1.3,
  scikit-learn 1.6.1, fastembed 0.4.2 &
  Estimation tarifaire, optimisation, prediction de pannes, assistant
  documentaire \\
Infrastructure & PostgreSQL 16, RabbitMQ 3.13, Qdrant 1.12, nginx 1.25
  et 1.27, Docker Compose &
  Persistance, courtier de messages, base vectorielle, passerelle et
  orchestration des vingt conteneurs \\
\bottomrule
\end{tabularx}
\end{table}

\begin{figure}[H]
  \centering
  \includegraphics[height=1.15cm,keepaspectratio]{logo-dotnet.png}\hspace{0.7cm}
  \includegraphics[height=1.15cm,keepaspectratio]{logo-csharp.png}\hspace{0.7cm}
  \includegraphics[height=1.15cm,keepaspectratio]{logo-react.png}\hspace{0.7cm}
  \includegraphics[height=1.15cm,keepaspectratio]{logo-typescript.png}\hspace{0.7cm}
  \includegraphics[height=1.15cm,keepaspectratio]{logo-python.jpg}\hspace{0.7cm}
  \includegraphics[height=1.15cm,keepaspectratio]{logo-postgres.png}\hspace{0.7cm}
  \includegraphics[height=1.15cm,keepaspectratio]{logo-rabbitmq.png}\hspace{0.7cm}
  \includegraphics[height=1.15cm,keepaspectratio]{logo-docker.png}
  \caption{Principales technologies mises en oeuvre}
  \label{fig:technologies}
\end{figure}

% =====================================================================
\section{Architecture logique}
\label{sec:archi-logique}

L'architecture logique organise les responsabilites en quatre niveaux
empiles, le flux descendant de l'interface vers la persistance.

\figReelle{logical-architecture.png}
  {Architecture logique de la plateforme, en quatre niveaux}
  {fig:archi-logique}{0.97}

\paragraph{Niveau 1, presentation.}
Trois applications servent trois publics : le back-office pour le
personnel, le portail pour les clients, l'application mobile pour les
chauffeurs. Aucune ne contient de regle metier ; elles consomment des API.

\paragraph{Niveau 2, passerelle.}
Un serveur nginx assure trois fonctions : le routage vers le service
concerne d'apres le prefixe de l'URL, la validation du jeton
d'authentification et le filtrage des origines autorisees. Il constitue
le point d'entree unique du systeme.

\paragraph{Niveau 3, coeur applicatif.}
Huit services metier en .NET et quatre services d'aide a la decision en
Python, relies par un courtier de messages. Une bibliotheque commune
porte les enumerations et les evenements d'integration partages.

\paragraph{Niveau 4, persistance.}
Une base relationnelle, une base vectorielle et un volume de fichiers.

\begin{definitionbox}{Une lecture a retenir}
Sur ce diagramme, les liaisons vont \exemple{de groupe a groupe} et non
de service a service. Les huit services metier partagent la meme
passerelle, le meme courtier et la meme base : les relier un a un
produirait une trentaine de traits sans rien apprendre de plus. Le detail
service par service figure dans les chapitres de sprint.
\end{definitionbox}

% =====================================================================
\section{Architecture physique}
\label{sec:archi-physique}

La vue physique decrit le deploiement effectif : vingt conteneurs sur un
hote unique, orchestres par Docker Compose.

\figReelle{physical-architecture.png}
  {Architecture physique : deploiement conteneurise}
  {fig:archi-physique}{0.97}

Trois points de cette vue meritent un commentaire.

\paragraph{La resolution des noms.}
A l'interieur du reseau, les services s'appellent par leur nom et non par
une adresse. Les ports 5001 a 5009 visibles sur le diagramme ne sont que
des publications vers l'hote, utiles au diagnostic et aux appels directs
du portail client.

\paragraph{L'ordre de demarrage.}
Il repose sur une condition de bonne sante : la base et le courtier
demarrent d'abord, puis les services metier, puis la passerelle. Les huit
services .NET exposent un point de controle de sante que
l'orchestrateur interroge.

\paragraph{Les acces directs.}
Le portail client et l'application chauffeur ne passent pas systematiquement
par la passerelle : le suivi de position et l'estimation tarifaire sont
atteints en direct. Ce n'est pas un oubli mais une consequence du
perimetre de routage de la passerelle, examinee en
section~\ref{sec:ecarts}.

% =====================================================================
\section{Choix structurants}
\label{sec:choix}

\subsection{Clean Architecture au sein de chaque service}

Chaque service metier est decoupe en quatre projets, dont les
dependances de compilation sont contraintes.

\figReelle{platform-clean-architecture.png}
  {Decoupage en couches d'un service metier}
  {fig:clean-archi}{0.52}

La regle est la suivante : \essentiel{les dependances pointent vers
l'interieur}. La couche \exemple{Domain}, qui porte les entites et les
regles de gestion, ne depend de rien d'autre que de la bibliotheque
commune. Elle ignore la base de donnees, le protocole HTTP et le
courtier de messages.

L'inversion se joue dans la couche \exemple{Infrastructure} : celle-ci
\essentiel{implemente} des interfaces declarees dans la couche
\exemple{Application}, et non l'inverse. Concretement, un service
applicatif manipule une interface de depot sans savoir qu'elle est
satisfaite par Entity Framework Core.

L'interet n'est pas theorique. Le domaine se teste sans base de donnees,
et le remplacement d'une brique technique n'affecte pas les regles
metier. Le service de maintenance predictive en fournit l'illustration :
son moteur de prediction est declare par une interface, et l'implantation
actuelle, fondee sur des regles deterministes, peut ceder la place a un
modele d'apprentissage sans qu'aucune autre couche ne bouge.

\subsection{Deux bus de messages plutot qu'un}

La plateforme utilise le meme courtier RabbitMQ de deux manieres
distinctes, ce qui surprend a la premiere lecture.

\begin{table}[H]
\centering
\caption{Les deux usages du courtier de messages}
\label{tab:deux-bus}
\small
\begin{tabularx}{\textwidth}{@{}p{3.2cm}XX@{}}
\toprule
 & \textbf{MassTransit} & \textbf{Echange par sujet} \\ \midrule
Format & Message type, avec enveloppe & JSON brut, sans enveloppe \\
Emetteurs & Services .NET & Services .NET \\
Destinataires & Services .NET & Services Python \\
Motif & Contrat fortement type entre services d'une meme pile &
  Les consommateurs Python liraient difficilement l'enveloppe
  MassTransit ; le JSON brut leur est directement exploitable \\
\bottomrule
\end{tabularx}
\end{table}

Ce choix est un compromis assume : il introduit deux conventions a
maintenir, mais evite d'imposer aux services Python le decodage d'un
format concu pour l'ecosysteme .NET.

\subsection{Ecarts assumes par rapport au patron de reference}
\label{sec:ecarts}

Une architecture microservices de reference prevoit une base de donnees
par service. \essentiel{Ce n'est pas le cas ici : les neuf chaines de
connexion pointent toutes vers la meme base.} Cet ecart doit etre
explicite plutot que masque.

\begin{table}[H]
\centering
\caption{Ecarts par rapport au patron de reference}
\label{tab:ecarts}
\small
\begin{tabularx}{\textwidth}{@{}p{3.6cm}X@{}}
\toprule
\textbf{Ecart} & \textbf{Portee et consequence} \\ \midrule
Base partagee &
  Le cloisonnement est LOGIQUE : chaque service ne possede en ecriture
  que ses propres tables et lit celles des autres domaines en SQL direct.
  C'est pourquoi certaines entites sont redeclarees en lecture seule dans
  plusieurs couches Domain. Consequence : le decouplage des donnees
  repose sur une discipline de developpement, non sur l'infrastructure. \\
Absence de cache &
  Aucun cache applicatif n'est deploye. La limitation de debit de
  l'assistant est tenue en memoire du processus, ce qui ne resisterait
  pas a une mise a l'echelle horizontale. \\
Absence de terminaison TLS &
  La passerelle publie en clair. Le chiffrement devra etre ajoute avant
  toute mise en ligne. \\
Hote unique &
  Le deploiement repose sur Docker Compose, sur une seule machine. Il n'y
  a ni orchestrateur de grappe ni redondance. \\
\bottomrule
\end{tabularx}
\end{table}

Ces ecarts sont le prix d'un projet mene en \stageDuree{} mois par une
seule personne. Ils sont documentes dans les diagrammes d'architecture
eux-memes, sous forme de notes, afin qu'un lecteur ne les prenne pas pour
des oublis.

% =====================================================================
\section*{Conclusion}
\addcontentsline{toc}{section}{Conclusion}

Ce chapitre a decrit les moyens et l'architecture. La plateforme repose
sur quatre niveaux logiques, deployes en vingt conteneurs sur un hote
unique. Chaque service metier applique le decoupage en couches de la
Clean Architecture, dont la regle de dependance protege le domaine des
choix techniques. Deux usages du courtier de messages coexistent
volontairement, pour ne pas imposer aux services Python un format concu
pour .NET. Quatre ecarts par rapport au patron de reference ont enfin ete
identifies et justifies.

Les six chapitres suivants presentent la realisation, sprint par sprint.
Le premier d'entre eux traite du socle sur lequel tout le reste
s'appuie : les identites et les acces.

% =====================================================================
%  Chapitre 4 - Sprint 1 : identites, acces et securite
% =====================================================================
\chapter{Sprint 1 : identites, acces et securite}
\label{chap:sprint1}

\section*{Introduction}
\addcontentsline{toc}{section}{Introduction}

Ce premier sprint construit le socle sur lequel reposent tous les
suivants : sans identites ni habilitations, aucun ecran de la plateforme
ne peut etre protege et aucune action ne peut etre imputee a son auteur.
Il couvre l'inscription, l'authentification, la gestion du compte
personnel et l'administration des comptes par les responsables.

% =====================================================================
\section{Planification du sprint}

\subsection{Objectifs}

\begin{enumerate}
  \item Permettre a un visiteur de creer un compte et d'en confirmer
    l'adresse electronique.
  \item Authentifier les utilisateurs et proteger les acces par un second
    facteur optionnel.
  \item Offrir au chauffeur un mode de connexion adapte a la mobilite.
  \item Donner a l'administrateur les moyens de statuer sur les demandes
    d'inscription et de gerer le cycle de vie des comptes.
  \item Rendre toute connexion tracable.
\end{enumerate}

\subsection{Carnet du sprint}

\begin{table}[H]
\centering
\caption{Carnet du sprint 1}
\label{tab:backlog-s1}
\small
\begin{tabularx}{\textwidth}{@{}cp{2.5cm}Xcc@{}}
\toprule
\textbf{N\textsuperscript{o}} & \textbf{Acteur} & \textbf{Besoin exprime} &
\textbf{Prio.} & \textbf{Pts} \\ \midrule
1 & Visiteur & Creer un compte & Must & 5 \\
2 & Visiteur & Confirmer son adresse electronique & Must & 3 \\
3 & Utilisateur & Se connecter a son espace & Must & 5 \\
4 & Utilisateur & Activer un second facteur & Should & 8 \\
5 & Chauffeur & Se connecter par code PIN & Should & 5 \\
6 & Utilisateur & Gerer son profil et son mot de passe & Must & 3 \\
7 & Administrateur & Statuer sur les demandes d'inscription & Must & 5 \\
8 & Responsable & Gerer les habilitations d'un compte & Must & 3 \\
9 & Administrateur & Archiver un compte sans perdre son historique & Must & 5 \\
10 & Administrateur & Consulter l'historique des connexions & Should & 3 \\
\midrule
\multicolumn{4}{@{}r}{\textbf{Total}} & \textbf{45} \\
\bottomrule
\end{tabularx}
\end{table}

% =====================================================================
\section{Elaboration technique}

\subsection{Modelisation fonctionnelle}

\figReelle{iam-usecase.png}
  {Cas d'utilisation du service des identites et des acces}
  {fig:uc-iam}{0.50}

Le diagramme de la figure~\ref{fig:uc-iam} fait apparaitre trois
regroupements. \essentiel{L'acces au service} rassemble ce qui precede
l'entree dans l'application, \essentiel{le compte personnel} ce que
l'utilisateur gere pour lui-meme, et \essentiel{l'administration des
comptes} ce qui est reserve au responsable et a l'administrateur.

Deux relations d'extension meritent d'etre relevees. \exemple{Valider un
second facteur} etend \exemple{se connecter} : l'etape n'a lieu que si
l'utilisateur a active la protection. \exemple{Archiver un compte} etend
\exemple{supprimer un compte} : l'archivage est propose lorsque le compte
porte un historique metier. Le cas \exemple{consulter les dependances
metier} est en revanche \essentiel{inclus}, la verification etant
systematique avant toute suppression.

\begin{table}[H]
\centering
\caption{Description textuelle du cas « se connecter »}
\label{tab:desc-connexion}
\small
\begin{tabularx}{\textwidth}{@{}p{2.5cm}X@{}}
\toprule
\textbf{Rubrique} & \textbf{Contenu} \\ \midrule
Acteur & Utilisateur \\
Precondition & Le compte existe et son statut est ACTIF \\
Scenario nominal & Saisie des identifiants, verification de l'empreinte,
  emission d'un jeton d'acces valable soixante minutes et d'un jeton de
  rafraichissement, journalisation de la connexion \\
Extension & Si le second facteur est actif, un jeton temporaire est emis
  et le code est demande avant la delivrance du jeton d'acces \\
Alternatives & Identifiants invalides : compteur d'echecs incremente,
  verrouillage de quinze minutes au cinquieme echec. Compte en attente,
  suspendu ou archive : refus motive \\
\bottomrule
\end{tabularx}
\end{table}

\subsection{Conception statique}

\figReelle{iam-class-domain.png}
  {Modele du domaine du service des identites et des acces}
  {fig:dom-iam}{0.52}

Le modele repose sur une \essentiel{hierarchie d'heritage}. La classe
\exemple{Utilisateur} porte ce qui est commun ; \exemple{Client} et
\exemple{Employe} la specialisent, et \exemple{Chauffeur} specialise a son
tour \exemple{Employe}. Cette hierarchie est stockee dans une table
unique, distinguee par un discriminant : les specialisations partagent le
meme identifiant, ce qui evite une jointure a chaque lecture.

Le \essentiel{statut} du compte est un attribut a part entiere, distinct
de l'indicateur d'activite. Il decrit un cycle de vie en cinq etats : en
attente, actif, suspendu, desactive, archive. \essentiel{L'archivage n'est
pas une suppression} : la ligne est conservee, ce qui preserve les
missions, les factures et les interventions qui la referencent, et
l'operation reste reversible.

\subsection{Conception dynamique}

\figReelle{iam-seq-inscription.png}
  {Sequence : inscription d'un compte}
  {fig:seq-inscription}{0.72}

La sequence de la figure~\ref{fig:seq-inscription} montre que le compte
est cree au statut \essentiel{en attente} : il reste inutilisable tant
qu'un administrateur ne l'a pas approuve. Deux effets accompagnent la
creation : la publication d'un evenement sur le bus et l'envoi d'un
courriel de confirmation valable vingt-quatre heures.

\figReelle{iam-seq-approbation-compte.png}
  {Sequence : approbation d'un compte par l'administrateur}
  {fig:seq-approbation}{0.72}

L'approbation, figure~\ref{fig:seq-approbation}, fait passer le compte a
l'etat actif et publie a son tour un evenement. La lecture des deux
diagrammes met en evidence une constante de conception :
\essentiel{le controleur ne s'adresse jamais directement au depot}. Il
passe par une interface de service, satisfaite par une implantation qui
seule connait la base. C'est la regle de dependance exposee au
chapitre~\ref{chap:architecture}.

\subsection{Integration avec les autres services}

\figReelle{iam-communication.png}
  {Echanges du service des identites avec le reste de la plateforme}
  {fig:com-iam}{0.62}

Le service des identites est le seul a \essentiel{publier} des evenements
types sur le bus de messages. Le service de messagerie les consomme afin
de creer automatiquement les conversations d'accueil. Ce couplage est
\essentiel{asynchrone} : si la messagerie est indisponible, l'inscription
aboutit malgre tout et l'evenement sera traite au redemarrage.

% =====================================================================
\section{Realisation et tests de validation}

\figTrois{s1-connexion.jpg}{s1-double-facteur.jpg}{s1-connexion-pin.jpg}
  {Parcours d'authentification : connexion, validation du second facteur,
   connexion par code PIN du chauffeur}
  {fig:ecrans-auth}{0.40}

\figDeux{s1-comptes-attente.jpg}{s1-comptes-archives.jpg}
  {Administration des comptes : demandes en attente et comptes archives
   restaurables}
  {fig:ecrans-admin}{0.55}

\begin{table}[H]
\centering
\caption{Tests de validation du sprint 1}
\label{tab:tests-s1}
\small
\begin{tabularx}{\textwidth}{@{}cXp{3.0cm}c@{}}
\toprule
\textbf{N\textsuperscript{o}} & \textbf{Cas teste} & \textbf{Attendu} &
\textbf{Etat} \\ \midrule
T1 & Identifiants valides & Jeton delivre, connexion journalisee & Reussi \\
T2 & Cinq echecs consecutifs & Verrouillage de quinze minutes & Reussi \\
T3 & Connexion sur un compte en attente & Refus motive & Reussi \\
T4 & Second facteur actif & Code demande avant le jeton & Reussi \\
T5 & Archivage d'un compte porteur de missions & Missions preservees & Reussi \\
T6 & Restauration d'un compte archive & Compte de nouveau actif & Reussi \\
\bottomrule
\end{tabularx}
\end{table}

% =====================================================================
\section*{Conclusion}
\addcontentsline{toc}{section}{Conclusion}

Ce sprint livre le socle d'identites de la plateforme : inscription
soumise a approbation, authentification protegee par un second facteur
optionnel, connexion par code PIN pour les chauffeurs, gestion des
habilitations et archivage reversible. Les quarante-cinq points du carnet
ont ete realises et les six tests de validation sont passes.

Le choix le plus structurant de ce sprint n'est pas technique mais
metier : \essentiel{ne jamais supprimer un compte qui porte un
historique}. Il conditionne tous les modules suivants, ou les missions,
les factures et les interventions referencent leur auteur.

Le sprint suivant construit le second socle, celui des ressources
materielles : le referentiel de la flotte et des equipements.

% =====================================================================
%  Chapitre 5 - Sprint 2 : referentiel de flotte et d'equipements
% =====================================================================
\chapter{Sprint 2 : referentiel de flotte et d'equipements}
\label{chap:sprint2}

\section*{Introduction}
\addcontentsline{toc}{section}{Introduction}

Le sprint precedent a etabli qui utilise la plateforme. Celui-ci etablit
\essentiel{avec quoi l'entreprise travaille} : le parc de vehicules, les
equipements embarques, les pieces administratives et le carburant. C'est
le second socle du systeme, celui sur lequel s'appuieront les missions et
la maintenance : on n'affecte pas un vehicule qui n'existe pas au
referentiel.

% =====================================================================
\section{Planification du sprint}

\subsection{Objectifs}

\begin{enumerate}
  \item Constituer un referentiel unique des vehicules et des
    equipements, chacun decrit une seule fois.
  \item Rattacher a chaque vehicule ses pieces administratives et ses
    assurances, avec leurs dates d'echeance.
  \item Suivre le cycle de vie d'un vehicule par son statut
    d'exploitation.
  \item Tracer l'affectation et la restitution des equipements.
  \item Enregistrer les pleins de carburant et en tirer une consommation.
\end{enumerate}

\subsection{Carnet du sprint}

\begin{table}[H]
\centering
\caption{Carnet du sprint 2}
\label{tab:backlog-s2}
\small
\begin{tabularx}{\textwidth}{@{}cp{2.5cm}Xcc@{}}
\toprule
\textbf{N\textsuperscript{o}} & \textbf{Acteur} & \textbf{Besoin exprime} &
\textbf{Prio.} & \textbf{Pts} \\ \midrule
1 & Utilisateur & Consulter le parc et rechercher un vehicule disponible & Must & 5 \\
2 & Responsable & Tenir la fiche d'un vehicule & Must & 8 \\
3 & Responsable & Enregistrer les pieces administratives et les assurances & Must & 5 \\
4 & Utilisateur & Consulter les echeances a venir & Must & 3 \\
5 & Responsable & Tenir la fiche d'un equipement & Must & 5 \\
6 & Utilisateur & Affecter et restituer un equipement & Must & 5 \\
7 & Utilisateur & Realiser une inspection d'equipement & Should & 3 \\
8 & Utilisateur & Declarer un plein de carburant & Should & 3 \\
9 & Utilisateur & Suivre la consommation de la flotte & Should & 5 \\
10 & Administrateur & Sortir un bien de l'inventaire & Must & 3 \\
\midrule
\multicolumn{4}{@{}r}{\textbf{Total}} & \textbf{45} \\
\bottomrule
\end{tabularx}
\end{table}

% =====================================================================
\section{Elaboration technique}

\subsection{Modelisation fonctionnelle}

\figReelle{fleet-usecase.png}
  {Cas d'utilisation du service de gestion de flotte}
  {fig:uc-fleet}{0.75}

La figure~\ref{fig:uc-fleet} distingue trois domaines. \essentiel{L'exploitation
courante} rassemble ce qui est ouvert a tout utilisateur authentifie : la
consultation du parc ne porte en effet aucune restriction de role dans le
code, seules les ecritures en sont assorties. \essentiel{La tenue du
referentiel} revient au responsable, et \essentiel{le retrait definitif}
au seul administrateur.

Cette gradation reproduit une regle simple : \essentiel{lire est ouvert,
ecrire est encadre, supprimer est reserve}. Elle se retrouvera dans tous
les modules suivants.

\begin{table}[H]
\centering
\caption{Description textuelle du cas « immatriculer un vehicule »}
\label{tab:desc-immat}
\small
\begin{tabularx}{\textwidth}{@{}p{2.5cm}X@{}}
\toprule
\textbf{Rubrique} & \textbf{Contenu} \\ \midrule
Acteur & Responsable \\
Precondition & Le vehicule n'existe pas encore au referentiel \\
Scenario nominal & Saisie de l'immatriculation, du numero de chassis, de
  la marque, du modele et des caracteristiques techniques ; controle
  d'unicite de l'immatriculation et du chassis ; enregistrement au statut
  DISPONIBLE \\
Alternatives & Immatriculation deja presente : creation refusee. Champ
  obligatoire manquant : le formulaire signale la ligne fautive \\
Postcondition & Le vehicule est visible au parc et peut recevoir des
  equipements, des documents et des missions \\
\bottomrule
\end{tabularx}
\end{table}

\subsection{Conception statique}

\figReelle{fleet-class-domain.png}
  {Modele du domaine du service de gestion de flotte}
  {fig:dom-fleet}{0.90}

Le modele s'organise autour de la classe \exemple{Vehicule}, a laquelle
se rattachent quatre familles d'objets : les \exemple{documents
administratifs} et les \exemple{assurances}, qui portent des dates
d'echeance ; les \exemple{equipements}, affectes puis restitues ; les
\exemple{pleins de carburant}, rattaches a une source ; et les
\exemple{inspections}.

Deux points de conception meritent d'etre releves.

\paragraph{L'echeance est une donnee, pas un calcul.}
La date de fin de validite est stockee sur le document et sur
l'assurance. Le service d'alertes la lit pour produire ses avertissements
plutot que de la deduire : une echeance depassee reste visible meme si
personne n'a ouvert l'ecran depuis des semaines.

\paragraph{L'equipement a son propre cycle de vie.}
Il porte un etat, un compteur d'utilisations et des seuils d'usure. Il
peut etre affecte a un vehicule ou rester au magasin : la relation est
donc optionnelle des deux cotes, ce que traduit la multiplicite du
diagramme.

\subsection{Conception dynamique}

% Ces deux sequences occupaient deja la meme largeur, la justification
% entiere. Elles paraissaient pourtant de tailles tres differentes, et
% elles l'etaient : la premiere mesure 1475 pixels de large, la seconde
% 895. A largeur egale, la premiere subit donc une reduction deux fois
% plus forte, et son texte devient deux fois plus petit.
%
% Egaliser les largeurs ne pouvait rien y faire. C'est le RAPPORT DE
% RENDU qu'il faut egaliser. La premiere passe en paysage : 24,7 cm au
% lieu de 16, soit un rendu a 0,63 de sa taille native. La seconde,
% laissee en pleine justification, rend a 0,68. Les deux textes ont
% desormais la meme taille a un vingtieme pres.
\figPaysage{fleet-seq-immatriculation.png}
  {Sequence : immatriculation d'un vehicule}
  {fig:seq-immat}

\figReelle{fleet-seq-suivi-echeances.png}
  {Sequence : consultation des echeances a venir}
  {fig:seq-echeances}{1.00}

La premiere sequence montre le controle d'unicite effectue avant tout
enregistrement. La seconde, figure~\ref{fig:seq-echeances}, illustre un
mecanisme different : \essentiel{aucune ecriture n'a lieu}, le service
agrege les dates portees par les documents, les assurances et les permis,
puis les ordonne par proximite. C'est une lecture transversale, et c'est
elle qui alimentera le module d'alertes du sprint~5.

Une troisieme sequence, le changement d'etat d'un vehicule, suit la meme
structure que l'immatriculation et n'est pas reproduite ici.

% =====================================================================
\section{Realisation}

\subsection{Le parc et l'inventaire}

La figure~\ref{fig:s2-vehicules} presente l'ecran de tete du module : le
parc, dans son ensemble. Chaque ligne porte l'immatriculation, la marque,
le modele, le kilometrage et surtout le \essentiel{statut
d'exploitation} --- disponible, en mission, en maintenance, hors service.
Ce statut est la donnee la plus consultee de tout le referentiel, parce
que c'est elle qui conditionne l'affectation : un vehicule en maintenance
ne doit pas apparaitre dans la liste des candidats a une mission. Les
filtres places en tete permettent de restreindre l'affichage a un statut,
a un type de vehicule ou a une plage de kilometrage.

\capture{s2-liste-vehicules.jpg}
  {Parc de vehicules : liste filtrable et statuts d'exploitation}
  {fig:s2-vehicules}

L'inventaire des equipements, figure~\ref{fig:s2-equipements}, obeit a la
meme logique de presentation mais repond a une question differente :
\exemple{ou se trouve ce materiel ?} La colonne d'affectation indique soit
un vehicule, soit le magasin. Un hayon, un groupe frigorifique ou un jeu
de sangles passe de l'un a l'autre au fil des besoins, et c'est ce
mouvement que l'ecran doit rendre lisible. Le compteur d'utilisations,
visible sur la droite, prepare la maintenance preventive du sprint~4 : au
dela d'un seuil, l'equipement doit etre controle.

\capture{s2-liste-equipements.jpg}
  {Inventaire des equipements et leur affectation courante}
  {fig:s2-equipements}

\subsection{La fiche d'un vehicule}

La fiche detaillee est organisee en onglets, pour eviter un formulaire
unique de plusieurs ecrans de haut. Le premier, figure~\ref{fig:s2-fiche1},
rassemble l'identite technique : chassis, energie, charge utile,
kilometrage, date de mise en circulation. Ces valeurs changent peu, mais
elles alimentent des decisions frequentes --- la charge utile determine
quelles marchandises le vehicule peut prendre, l'energie sert au calcul de
consommation.

\capture{s2-fiche-vehicule-1.jpg}
  {Fiche d'un vehicule : identite et caracteristiques techniques}
  {fig:s2-fiche1}

Le second onglet, figure~\ref{fig:s2-fiche2}, porte les pieces
administratives : carte grise, visite technique, vignette, autorisations
de transport. Chacune est datee et possede une \essentiel{fin de
validite}. C'est ici que se materialise le choix de conception expose plus
haut : la date est saisie une fois, stockee, puis relue par le module
d'alertes. L'ecran signale par une pastille de couleur les pieces dont
l'echeance approche, ce qui donne au responsable une lecture immediate
sans qu'il ait a comparer les dates lui-meme.

\capture{s2-fiche-vehicule-2.jpg}
  {Fiche d'un vehicule : pieces administratives et echeances}
  {fig:s2-fiche2}

Enfin, la figure~\ref{fig:s2-detail-equip} montre le detail d'un
equipement affecte. On y retrouve son etat, son compteur d'utilisations,
le vehicule porteur et l'historique de ses affectations successives. Cet
historique n'est pas un ornement : lorsqu'un equipement tombe en panne, il
permet de savoir sur quels vehicules il a servi et donc quelles missions
ont pu en patir.

\capture{s2-detail-equipement.jpg}
  {Detail d'un equipement et historique de ses affectations}
  {fig:s2-detail-equip}

% =====================================================================
\section{Tests de validation}

\begin{table}[H]
\centering
\caption{Tests de validation du sprint 2}
\label{tab:tests-s2}
\small
\begin{tabularx}{\textwidth}{@{}cXp{3.0cm}c@{}}
\toprule
\textbf{N\textsuperscript{o}} & \textbf{Cas teste} & \textbf{Attendu} &
\textbf{Etat} \\ \midrule
T1 & Immatriculation d'un vehicule inedit & Vehicule cree, statut DISPONIBLE & Reussi \\
T2 & Immatriculation deja presente & Creation refusee & Reussi \\
T3 & Document dont l'echeance est depassee & Vehicule signale dans les echeances & Reussi \\
T4 & Affectation d'un equipement a un vehicule & Equipement indisponible au magasin & Reussi \\
T5 & Restitution d'un equipement & Compteur d'utilisations incremente & Reussi \\
T6 & Plein de carburant enregistre & Consommation moyenne recalculee & Reussi \\
\bottomrule
\end{tabularx}
\end{table}

% =====================================================================
\section*{Conclusion}
\addcontentsline{toc}{section}{Conclusion}

Ce sprint livre le referentiel materiel de la plateforme : seize
vehicules decrits avec leurs caracteristiques, leurs pieces
administratives et leurs assurances, quatorze equipements suivis avec
leur etat et leur affectation, et le suivi du carburant. Les
quarante-cinq points du carnet ont ete realises.

L'apport le plus durable de ce sprint n'est pas un ecran mais une
donnee : \essentiel{les dates d'echeance sont desormais stockees et
interrogeables}. Sans elles, la surveillance reglementaire du sprint~5
n'aurait aucune matiere sur laquelle travailler.

Les deux socles etant poses, le sprint suivant aborde le coeur de
l'activite : la chaine qui va de la demande du client a la livraison de
sa marchandise.

% =====================================================================
%  Chapitre 6 - Sprint 3 : demande, devis, mission et suivi en direct
% =====================================================================
\chapter{Sprint 3 : demande, devis, mission et suivi en direct}
\label{chap:sprint3}

\section*{Introduction}
\addcontentsline{toc}{section}{Introduction}

Les deux sprints precedents ont pose les socles : qui utilise la
plateforme, et avec quelles ressources. Celui-ci aborde \essentiel{le
coeur de l'activite} : la chaine qui va de la demande du client a la
livraison de sa marchandise, en passant par le devis, l'affectation et le
suivi. C'est le sprint le plus dense du projet, et il mobilise deux
services : celui des missions et le moteur d'estimation tarifaire.

% =====================================================================
\section{Planification du sprint}

\subsection{Objectifs}

\begin{enumerate}
  \item Ouvrir au client un canal autonome de demande de transport.
  \item Estimer le prix d'un trajet avant meme que le devis ne soit
    etabli.
  \item Formaliser le cycle du devis, de l'envoi a l'acceptation ou au
    refus.
  \item Creer la mission, l'affecter et en suivre l'execution.
  \item Restituer la position du vehicule en direct sur une carte.
\end{enumerate}

\subsection{Carnet du sprint}

\begin{table}[H]
\centering
\caption{Carnet du sprint 3}
\label{tab:backlog-s3}
\small
\begin{tabularx}{\textwidth}{@{}cp{2.3cm}Xcc@{}}
\toprule
\textbf{N\textsuperscript{o}} & \textbf{Acteur} & \textbf{Besoin exprime} &
\textbf{Prio.} & \textbf{Pts} \\ \midrule
1 & Client & Deposer une demande de transport & Must & 8 \\
2 & Client & Connaitre une fourchette de prix avant de valider & Should & 8 \\
3 & Responsable & Etablir et envoyer un devis & Must & 8 \\
4 & Client & Accepter ou refuser un devis recu & Must & 5 \\
5 & Responsable & Affecter un vehicule et un chauffeur & Must & 5 \\
6 & Chauffeur & Transmettre sa position depuis l'application mobile & Should & 8 \\
7 & Client & Suivre sa mission en direct sur une carte & Should & 8 \\
8 & Responsable & Cloturer ou annuler une mission & Must & 3 \\
9 & Client & Ouvrir un ticket de support sur une mission & Could & 3 \\
10 & Responsable & Consulter l'historique d'un trajet & Could & 3 \\
\midrule
\multicolumn{4}{@{}r}{\textbf{Total}} & \textbf{59} \\
\bottomrule
\end{tabularx}
\end{table}

% =====================================================================
\section{Elaboration technique}

\subsection{Modelisation fonctionnelle}

\figReelle{mission-usecase.png}
  {Cas d'utilisation du service des missions}
  {fig:uc-mission}{0.54}

Le diagramme de la figure~\ref{fig:uc-mission} organise le service en
quatre domaines : la demande et le devis, l'execution des missions, le
suivi en direct et l'assistance. Trois cas se detachent a droite de la
frontiere, atteints par inclusion ou par extension : l'estimation du
prix, la creation de la mission et l'affectation des ressources.

\paragraph{Une regle metier structurante.}
\essentiel{Deposer une demande n'est pas commander une mission.} La
demande de transport est un \exemple{devis} : elle vit sa propre vie,
recoit un prix, puis est acceptee ou refusee. \essentiel{La mission n'est
creee qu'a l'acceptation}, ce que traduit la relation d'inclusion entre
\exemple{accepter un devis} et \exemple{creer une mission}.

Cette regle a une consequence concrete sur l'interface de programmation :
le point d'entree qui permettait a un client de creer directement une
mission a ete \essentiel{retire}. Le laisser aurait ouvert un chemin
parallele contournant le devis, et donc la tarification.

\begin{table}[H]
\centering
\caption{Description textuelle du cas « accepter un devis »}
\label{tab:desc-devis}
\small
\begin{tabularx}{\textwidth}{@{}p{2.5cm}X@{}}
\toprule
\textbf{Rubrique} & \textbf{Contenu} \\ \midrule
Acteur & Client \\
Precondition & La demande porte le statut QUOTATION\_SENT et un prix
  final a ete fixe \\
Scenario nominal & Le client consulte le devis, l'accepte ; le systeme
  cree la mission au statut EN\_ATTENTE, rattache la demande a la mission,
  fait passer la demande au statut ACCEPTED et publie un evenement sur
  le bus \\
Alternatives & Refus : la demande passe au statut REJECTED avec le motif
  saisi, aucune mission n'est creee. Devis deja traite : l'action est
  refusee \\
Postcondition & La mission existe et attend l'affectation d'un vehicule
  et d'un chauffeur \\
\bottomrule
\end{tabularx}
\end{table}

\subsection{Conception statique}

\figReelle{mission-class-domain.png}
  {Modele du domaine du service des missions}
  {fig:dom-mission}{0.30}

Le modele distingue nettement \exemple{DemandeTransport}, qui porte la
negociation (villes, marchandise, prix estime, fourchette, prix final,
methode de tarification), et \exemple{Mission}, qui porte l'execution
(vehicule, chauffeur, kilometrages, cout reel). Le lien entre les deux
est \essentiel{optionnel et oriente} : une demande peut exister sans
mission, jamais l'inverse.

\paragraph{Trois methodes de tarification.}
Le prix final s'accompagne d'une methode indiquant comment il a ete
obtenu : propose par le moteur et retenu tel quel, modifie, ou fixe
manuellement. \essentiel{Cette information est conservee}, ce qui permet
de mesurer a posteriori dans quelle mesure l'entreprise suit les
recommandations du modele.

\subsection{Conception dynamique}

\figReelle{mission-seq-demande-transport.png}
  {Sequence : depot d'une demande de transport}
  {fig:seq-demande}{0.54}

\figReelle{mission-seq-devis.png}
  {Sequence : etablissement puis acceptation ou refus d'un devis}
  {fig:seq-devis}{0.54}

La premiere sequence montre l'appel au moteur d'estimation \essentiel{avant}
l'enregistrement : le client voit la fourchette tarifaire au moment ou il
saisit sa demande. La seconde, figure~\ref{fig:seq-devis}, comporte un
bloc alternatif : l'acceptation cree la mission et publie un evenement,
le refus se contente d'enregistrer le motif.

\figReelle{mission-seq-suivi-direct.png}
  {Sequence : suivi d'une mission en direct}
  {fig:seq-suivi}{0.54}

Le suivi, figure~\ref{fig:seq-suivi}, repose sur une boucle : tant que la
mission est en cours, l'application du chauffeur transmet sa position, que
le client consulte. \essentiel{Ces positions ne transitent pas par la
passerelle} : l'application mobile s'adresse directement au service des
missions. Le detail figure a la section~\ref{sec:archi-physique}.

\subsection{Integration avec les services d'aide a la decision}

Le moteur d'estimation est un service Python autonome. Il charge en
memoire un modele de gradient boosting entraine sur un historique de
trajets, ainsi qu'une matrice de distances entre les villes desservies.
Sur reception d'une demande, il calcule la distance, applique le modele et
renvoie un prix encadre d'une fourchette basse et haute.

\figReelle{mission-communication.png}
  {Echanges du service des missions}
  {fig:com-mission}{0.60}

La figure~\ref{fig:com-mission} montre les deux natures d'echange du
service. L'appel au moteur d'estimation est \essentiel{synchrone} : le
client attend le prix pour decider. La demande d'optimisation d'une
tournee est \essentiel{asynchrone}, publiee sur le bus et traitee en
differe. Le critere est simple : \essentiel{on attend ce dont on a besoin
sur le champ, on delegue le reste}.

% =====================================================================
\section{Realisation et tests de validation}

\figDeux{s3-portail-reservation-1.jpg}{s3-estimation-tarifaire-1.jpg}
  {Portail client : formulaire de demande et fourchette tarifaire
   proposee}
  {fig:ecrans-demande}{0.46}

\figTrois{s3-demandes-transport.jpg}{s3-suivi-devis-client.jpg}{s3-carte-suivi.jpg}
  {Traitement d'une demande cote back-office, suivi du devis cote client,
   et carte de suivi d'une mission en cours}
  {fig:ecrans-mission}{0.40}

\begin{table}[H]
\centering
\caption{Tests de validation du sprint 3}
\label{tab:tests-s3}
\small
\begin{tabularx}{\textwidth}{@{}cXp{2.9cm}c@{}}
\toprule
\textbf{N\textsuperscript{o}} & \textbf{Cas teste} & \textbf{Attendu} &
\textbf{Etat} \\ \midrule
T1 & Depot d'une demande & Statut PENDING\_QUOTATION, prix estime affiche & Reussi \\
T2 & Envoi d'un devis modifie & Methode de tarification MODIFIED conservee & Reussi \\
T3 & Acceptation d'un devis & Mission creee, demande rattachee & Reussi \\
T4 & Refus d'un devis & Motif enregistre, aucune mission creee & Reussi \\
T5 & Position transmise par le chauffeur & Point visible sur la carte du client & Reussi \\
T6 & Trajet inconnu de la matrice & Estimation refusee, saisie manuelle possible & Reussi \\
\bottomrule
\end{tabularx}
\end{table}

% =====================================================================
\section*{Conclusion}
\addcontentsline{toc}{section}{Conclusion}

Ce sprint livre la chaine commerciale complete, de la demande au suivi en
direct. Les cinquante-neuf points du carnet ont ete realises, ce qui en
fait le sprint le plus lourd du projet.

Deux decisions en ressortent. La premiere est metier : \essentiel{la
reservation est un devis, la mission n'existe qu'apres acceptation} ;
elle a conduit a retirer un point d'entree qui permettait de la
contourner. La seconde est une regle d'integration : \essentiel{synchrone
pour ce dont l'utilisateur a besoin sur le champ, asynchrone pour le
reste}, reprise telle quelle par les sprints suivants.

Le sprint suivant traite du maintien en condition operationnelle de la
flotte : maintenance curative, puis anticipation des pannes.

% % =====================================================================
%  Chapitre 7 - Sprint 4 : maintenance curative et predictive
% =====================================================================
\chapter{Sprint 4 : maintenance curative et predictive}
\label{chap:sprint4}

\section*{Introduction}
\addcontentsline{toc}{section}{Introduction}

Une flotte ne vaut que par sa disponibilite. Ce sprint traite du maintien
en condition operationnelle des vehicules, en deux temps : d'abord la
maintenance curative, qui organise le travail de l'atelier, puis
l'anticipation des pannes, qui vise a intervenir avant la defaillance
plutot qu'apres.

% =====================================================================
\section{Planification du sprint}

\subsection{Objectifs}

\begin{enumerate}
  \item Organiser le travail de l'atelier autour d'ordres de travail
    suivis par etat.
  \item Tracer les interventions et leur consommation de pieces.
  \item Tenir le stock de pieces de rechange et alerter sur les seuils
    bas.
  \item Conserver l'historique des changements d'etat.
  \item Estimer le risque de panne par vehicule et designer l'organe
    concerne.
\end{enumerate}

\subsection{Carnet du sprint}

\begin{table}[H]
\centering
\caption{Carnet du sprint 4}
\label{tab:backlog-s4}
\small
\begin{tabularx}{\textwidth}{@{}cp{2.3cm}Xcc@{}}
\toprule
\textbf{N\textsuperscript{o}} & \textbf{Acteur} & \textbf{Besoin exprime} &
\textbf{Prio.} & \textbf{Pts} \\ \midrule
1 & Responsable & Ouvrir et instruire un ordre de travail & Must & 8 \\
2 & Technicien & Faire avancer un ordre de travail par etat & Must & 5 \\
3 & Utilisateur & Consulter l'historique des changements d'etat & Should & 3 \\
4 & Technicien & Declarer et suivre une intervention & Must & 8 \\
5 & Technicien & Consommer ou restituer une piece de rechange & Must & 5 \\
6 & Responsable & Tenir le catalogue et reapprovisionner le stock & Must & 5 \\
7 & Technicien & Documenter un ordre de travail & Should & 3 \\
8 & Responsable & Consulter le risque de panne de la flotte & Could & 8 \\
9 & Technicien & Connaitre l'organe a risque d'un vehicule & Could & 5 \\
10 & Administrateur & Purger le dossier de maintenance & Must & 3 \\
\midrule
\multicolumn{4}{@{}r}{\textbf{Total}} & \textbf{53} \\
\bottomrule
\end{tabularx}
\end{table}

% =====================================================================
\section{Elaboration technique}

\subsection{Modelisation fonctionnelle}

\figReelle{maintenance-usecase.png}
  {Cas d'utilisation du service de maintenance}
  {fig:uc-maintenance}{0.60}

Le diagramme de la figure~\ref{fig:uc-maintenance} repartit le service en
cinq domaines. La lecture des associations revele une
\essentiel{repartition des roles conforme au metier} : le responsable
ouvre et instruit les ordres de travail, le technicien les fait avancer
et realise les interventions, l'administrateur seul purge le dossier.

Trois cas figurent hors des domaines, atteints par extension :
l'historique des statuts, les pieces en stock faible et la consommation
de pieces. \essentiel{Aucun n'est un but en soi} : on ne se connecte pas
pour consulter un historique, on y accede depuis l'ordre de travail que
l'on examine.

\begin{table}[H]
\centering
\caption{Description textuelle du cas « declarer une intervention »}
\label{tab:desc-intervention}
\small
\begin{tabularx}{\textwidth}{@{}p{2.5cm}X@{}}
\toprule
\textbf{Rubrique} & \textbf{Contenu} \\ \midrule
Acteur & Technicien \\
Precondition & Un ordre de travail existe et n'est pas cloture \\
Scenario nominal & Saisie du type d'intervention, de la description et du
  kilometrage ; enregistrement au statut EN\_ATTENTE ; rattachement a
  l'ordre de travail \\
Extension & Consommation de pieces : chaque piece prelevee decremente le
  stock et alimente le cout de l'intervention \\
Postcondition & L'intervention est visible dans le dossier du vehicule et
  son cout s'agrege a celui de l'ordre de travail \\
\bottomrule
\end{tabularx}
\end{table}

\subsection{Conception statique}

\figReelle{maintenance-class-domain.png}
  {Modele du domaine du service de maintenance}
  {fig:dom-maintenance}{1.00}

Le modele distingue deux notions que l'usage confond souvent.
\essentiel{L'ordre de travail} est une decision de gestion : il ouvre un
dossier, porte une priorite et un etat. \essentiel{L'intervention} est un
acte technique : elle a une duree, un cout et consomme des pieces. Un
ordre de travail peut donner lieu a plusieurs interventions.

Deux mecanismes de tracabilite le completent. \exemple{HistoriqueStatutOT}
conserve chaque changement d'etat avec son auteur et son horodatage.
\exemple{InterventionPiece} porte la quantite et le prix unitaire
\essentiel{au moment du prelevement, et non le prix courant} : le cout
d'une intervention passee ne change pas si le tarif du fournisseur
evolue.

\subsection{Conception dynamique}

\figReelle{maintenance-seq-ordre-travail.png}
  {Sequence : ouverture et avancement d'un ordre de travail}
  {fig:seq-ot}{1.00}

\figReelle{maintenance-seq-intervention.png}
  {Sequence : declaration d'une intervention avec consommation de pieces}
  {fig:seq-intervention}{1.00}

La premiere sequence montre l'ecriture systematique d'une ligne
d'historique a chaque changement d'etat. La seconde,
figure~\ref{fig:seq-intervention}, illustre l'effet de bord de la
consommation : le stock de la piece est decremente dans la meme unite de
travail que l'enregistrement de l'intervention. \essentiel{Les deux
reussissent ou echouent ensemble}, ce qui interdit qu'une piece soit
comptee consommee sans que l'intervention existe.

\subsection{Anticipation des pannes}

Le service de maintenance predictive est le huitieme et dernier service
metier ajoute au projet. Il illustre a lui seul l'exigence d'evolutivite
formulee au chapitre~\ref{chap:analyse} : \essentiel{son ajout n'a impose
de modifier aucun service existant}.

\figReelle{predictionpanne-usecase.png}
  {Cas d'utilisation du service de maintenance predictive}
  {fig:uc-prediction}{1.00}

\figReelle{predictionpanne-sequence.png}
  {Sequence : calcul du risque de panne d'un vehicule}
  {fig:seq-prediction}{1.00}

\paragraph{Ce que fait reellement le moteur.}
Il faut ici etre precis, car le terme \exemple{prediction} prete a
confusion. \essentiel{Le moteur actuellement en service n'est pas un
modele appris : c'est un ensemble d'heuristiques de maintenance.} Il
croise quatre facteurs calcules sur les donnees reelles de la
plateforme : l'usure kilometrique par organe, le vieillissement du
vehicule, la recurrence des pannes deja constatees et le retard de
revision.

Ce choix est assume et repose sur deux arguments. D'abord, le module est
\essentiel{utile immediatement}, sans jeu d'entrainement historique, que
l'entreprise ne possede pas encore. Ensuite, et surtout, chaque
prediction est \essentiel{explicable} : le moteur designe l'organe
concerne et la raison du signalement. Or c'est precisement l'explication
qu'un exploitant attend avant d'immobiliser un vehicule.

Le moteur est declare derriere une interface, et l'enumeration de la
source d'une prediction prevoit deja la valeur correspondant a un modele
appris. \essentiel{Le remplacement se fera en une ligne d'enregistrement
de dependance}, sans toucher aux couches superieures.

% =====================================================================
\section{Realisation}

\subsection{Le pilotage de l'atelier}

Les ordres de travail sont presentes en colonnes, figure~\ref{fig:s4-ot},
une colonne par etat d'avancement : a planifier, planifie, en cours,
termine. Cette disposition n'a pas ete choisie pour l'effet visuel. Elle
repond a la maniere dont le chef d'atelier raisonne : il ne cherche pas un
ordre precis, il cherche \exemple{ce qui est en attente} et \exemple{ce
qui bloque}. Une liste triee par date ne repond pas a cette question, un
tableau par etat si. Le deplacement d'une carte d'une colonne a l'autre
declenche le changement d'etat cote serveur et \essentiel{ecrit une ligne
d'historique}, conformement a la sequence commentee plus haut.

\capture{s4-ordres-travail.jpg}
  {Ordres de travail presentes par etat d'avancement}
  {fig:s4-ot}

L'intervention, figure~\ref{fig:s4-interventions}, est le travail
reellement effectue au titre d'un ordre. Elle porte le technicien, la
duree, les pieces consommees et le cout. La distinction entre l'ordre et
l'intervention n'est pas une subtilite de modelisation : un meme ordre
peut donner lieu a plusieurs interventions successives, quand une
reparation est reprise ou completee. Confondre les deux aurait rendu
impossible le calcul du cout reel d'immobilisation d'un vehicule.

\capture{s4-interventions.jpg}
  {Interventions realisees, avec leur technicien et leur cout}
  {fig:s4-interventions}

\subsection{Le stock et l'anticipation}

Le magasin de pieces, figure~\ref{fig:s4-pieces}, affiche pour chaque
reference la quantite disponible et le \essentiel{seuil d'alerte} en
dessous duquel elle doit etre reapprovisionnee. Le stock n'est jamais
saisi a la main apres une reparation : il est decremente par
l'enregistrement de l'intervention, dans la meme unite de travail. Cette
mecanique, invisible a l'ecran, est ce qui garantit que la quantite
affichee corresponde a la realite de l'etagere.

\capture{s4-pieces-rechange.jpg}
  {Stock de pieces de rechange et seuils de reapprovisionnement}
  {fig:s4-pieces}

Le tableau de bord des risques, figure~\ref{fig:s4-prediction}, restitue
le travail du moteur decrit a la section precedente. Chaque vehicule y
porte un niveau de risque, l'organe concerne et \essentiel{la raison du
signalement} : depassement du kilometrage de revision, recurrence d'une
panne, anciennete. C'est cette derniere colonne qui fait la valeur de
l'ecran. Un score seul serait recu avec scepticisme par un exploitant qui
connait ses vehicules ; un score assorti de son motif se discute, et se
verifie.

\capture{s4-prediction-pannes.jpg}
  {Tableau de bord des risques de panne, par vehicule et par organe}
  {fig:s4-prediction}

% =====================================================================
\section{Tests de validation}

\begin{table}[H]
\centering
\caption{Tests de validation du sprint 4}
\label{tab:tests-s4}
\small
\begin{tabularx}{\textwidth}{@{}cXp{2.9cm}c@{}}
\toprule
\textbf{N\textsuperscript{o}} & \textbf{Cas teste} & \textbf{Attendu} &
\textbf{Etat} \\ \midrule
T1 & Changement d'etat d'un ordre de travail & Ligne d'historique ecrite & Reussi \\
T2 & Consommation de deux pieces & Stock decremente, cout agrege & Reussi \\
T3 & Piece descendue sous son seuil & Alerte de stock faible levee & Reussi \\
T4 & Prix fournisseur modifie apres coup & Cout de l'intervention inchange & Reussi \\
T5 & Vehicule au-dela du seuil de revision & Risque signale, organe designe & Reussi \\
T6 & Purge du dossier par l'administrateur & Elements retires, vehicule conserve & Reussi \\
\bottomrule
\end{tabularx}
\end{table}

% =====================================================================
\section*{Conclusion}
\addcontentsline{toc}{section}{Conclusion}

Ce sprint livre l'atelier complet : ordres de travail suivis par etat,
interventions tracees, stock surveille et tableau de bord des risques.
Les cinquante-trois points du carnet ont ete realises.

Deux principes le structurent. Le premier est comptable : \essentiel{le
prix d'une piece est fige au prelevement}, faute de quoi l'historique des
couts se reecrirait a chaque revision tarifaire. Le second est
methodologique : \essentiel{ce qui est presente comme une prediction doit
dire de quoi il retourne}. Un moteur de regles explicables a ete prefere
a un modele appris que les donnees disponibles n'auraient pas permis
d'entrainer serieusement, et l'interface qui les separe rend la bascule
possible sans reprise.

Le sprint suivant traite de ce que la maintenance et les missions
produisent en aval : la facturation, les alertes metier et le journal des
evenements.

% % =====================================================================
%  Chapitre 8 - Sprint 5 : facturation, alertes et journal metier
% =====================================================================
\chapter{Sprint 5 : facturation, alertes et journal metier}
\label{chap:sprint5}

\section*{Introduction}
\addcontentsline{toc}{section}{Introduction}

Les quatre sprints precedents ont produit de la matiere : des comptes,
des vehicules, des missions, des interventions. Ce sprint traite de trois
services qui n'en creent aucune : ils \essentiel{exploitent ce que les
autres ont produit}. La facturation transforme les missions en creances,
les alertes transforment les donnees en signaux, le journal metier
transforme les evenements en information de pilotage.

Cette position d'aval explique qu'ils aient ete traites ensemble, et
tardivement : chacun suppose que les modules amont existent deja.

% =====================================================================
\section{Planification du sprint}

\subsection{Objectifs}

\begin{enumerate}
  \item Etablir des factures a partir des missions realisees, sans
    double facturation possible.
  \item Suivre le cycle de vie d'une facture et identifier les impayes.
  \item Surveiller automatiquement les echeances, les stocks et les
    retards, et en alerter les responsables.
  \item Consigner les evenements metier dans un journal consultable.
  \item Diffuser les alertes critiques hors de l'application.
\end{enumerate}

\subsection{Carnet du sprint}

\begin{table}[H]
\centering
\caption{Carnet du sprint 5}
\label{tab:backlog-s5}
\small
\begin{tabularx}{\textwidth}{@{}cp{2.3cm}Xcc@{}}
\toprule
\textbf{N\textsuperscript{o}} & \textbf{Acteur} & \textbf{Besoin exprime} &
\textbf{Prio.} & \textbf{Pts} \\ \midrule
1 & Responsable & Etablir une facture depuis les missions facturables & Must & 8 \\
2 & Responsable & Suivre le cycle de vie d'une facture & Must & 5 \\
3 & Responsable & Identifier les factures impayees & Should & 3 \\
4 & Client & Consulter ses propres factures & Must & 3 \\
5 & Responsable & Consulter et filtrer les alertes metier & Must & 5 \\
6 & Responsable & Lancer les controles metier & Should & 5 \\
7 & Administrateur & Supprimer des alertes traitees & Should & 3 \\
8 & Responsable & Consulter le journal des evenements metier & Could & 5 \\
9 & Responsable & Recevoir une alerte critique hors application & Could & 5 \\
10 & Responsable & Editer un rapport quotidien de synthese & Could & 5 \\
\midrule
\multicolumn{4}{@{}r}{\textbf{Total}} & \textbf{47} \\
\bottomrule
\end{tabularx}
\end{table}

% =====================================================================
\section{Facturation}

\figReelle{finance-usecase.png}
  {Cas d'utilisation du service de facturation}
  {fig:uc-finance}{0.53}

\figReelle{finance-sequence.png}
  {Sequence : etablissement d'une facture}
  {fig:seq-finance}{0.61}

Le cas \exemple{etablir une facture} \essentiel{inclut} la verification
des missions facturables : la liste proposee au gestionnaire ne contient
que les missions terminees qui ne sont pas deja rattachees a une facture.
\essentiel{La double facturation est donc rendue impossible en amont},
plutot que detectee apres coup.

Le montant n'est pas saisi : il est calcule a partir du cout de la
mission, auquel s'ajoute la taxe sur la valeur ajoutee au taux en
vigueur. Les trois montants sont conserves separement, hors taxe, taxe et
toutes taxes comprises, ce qui evite de recalculer un historique si le
taux change.

Une facture parcourt cinq etats : brouillon, emise, payee, annulee et en
retard. Le passage en retard n'est pas une action de l'utilisateur mais
\essentiel{une consequence du temps} : il est deduit de la date
d'echeance, et c'est le service d'alertes qui le signale.

% =====================================================================
\section{Alertes metier}

\figReelle{alerts-usecase.png}
  {Cas d'utilisation du service d'alertes}
  {fig:uc-alerts}{0.69}

\figReelle{alerts-sequence.png}
  {Sequence : execution des controles metier}
  {fig:seq-alerts}{0.55}

Le service d'alertes ne detient aucune donnee propre : \essentiel{il lit
celles des autres domaines} et en tire des signaux. La
figure~\ref{fig:seq-alerts} montre le mecanisme : un balayage parcourt
les echeances documentaires, les assurances, les permis, les seuils de
stock et les missions en retard, puis produit une alerte pour chaque
situation rencontree.

\paragraph{Deux precautions de conception.}
La premiere est l'\essentiel{ecart des doublons} : le controle verifie
qu'une alerte identique et encore active n'existe pas deja avant d'en
creer une nouvelle. Sans cela, un vehicule dont la visite technique est
expiree produirait une alerte a chaque execution.

La seconde tient au \essentiel{vocabulaire}. Le type d'une alerte n'est
pas une chaine libre mais une valeur d'enumeration, propre au service.
Cette enumeration compte quarante et une valeurs et \essentiel{differe de
celle du projet partage} : ecrire une valeur du projet partage qui n'y
figure pas fait echouer la lecture. La contrainte a ete rencontree en
conditions reelles lors de la constitution du jeu de demonstration.

% =====================================================================
\section{Journal metier}

\figReelle{iboa-usecase.png}
  {Cas d'utilisation du service de journal metier}
  {fig:uc-iboa}{0.84}

\figReelle{iboa-communication.png}
  {Echanges du service de journal metier}
  {fig:com-iboa}{0.66}

Le service de journal metier occupe une position particuliere :
\essentiel{il ne consulte rien, il ecoute}. Ecrit en Python, il souscrit
a l'echange par sujet du courtier de messages et recoit les evenements
publies par les services metier. Pour chacun, il evalue une criticite,
regroupe les evenements lies et, au-dela d'un seuil, diffuse une alerte
par courriel ou par message instantane.

C'est le seul service de la plateforme dont \essentiel{le point d'entree
principal n'est pas une requete mais un message}. Son interface web ne
sert qu'a consulter ce qu'il a deja consigne.

\begin{table}[H]
\centering
\caption{Entites principales des trois services}
\label{tab:entites-s5}
\small
\begin{tabularx}{\textwidth}{@{}p{2.6cm}p{3.4cm}X@{}}
\toprule
\textbf{Service} & \textbf{Entite} & \textbf{Role et particularite} \\ \midrule
Facturation & Facture & Montants hors taxe, taxe et toutes taxes
  conserves separement ; rattachement optionnel a une mission \\
Facturation & Devis & Montants et validite ; distinct de la demande de
  transport du sprint~3, qu'il prolonge cote comptable \\
Alertes & Alerte & Type, priorite, statut, et references facultatives
  vers le vehicule, le chauffeur, l'equipement ou la mission concernes \\
Journal metier & Evenement & Charge utile brute recue du courtier,
  criticite evaluee, horodatage de reception \\
Journal metier & Alerte metier & Produite par regroupement d'evenements
  lies ; porte le canal de diffusion utilise \\
Journal metier & Preferences & Seuils et canaux choisis par le
  responsable pour la diffusion \\
\bottomrule
\end{tabularx}
\end{table}

% =====================================================================
\section{Realisation}

\subsection{La facturation}

La liste des factures, figure~\ref{fig:s5-factures}, est organisee autour
du \essentiel{statut de reglement} : brouillon, emise, payee, en retard,
annulee. C'est la seule lecture qui interesse reellement le service
administratif, dont la question quotidienne n'est pas \exemple{combien
avons-nous facture} mais \exemple{qu'est-ce qui n'est pas encore rentre}.
Les montants hors taxe, la taxe et le montant toutes taxes comprises sont
affiches separement, conformement au choix de modelisation decrit plus
haut : les recalculer a l'affichage aurait expose le document a des
divergences d'arrondi entre l'ecran et le document remis au client.

\capture{s5-liste-factures.jpg}
  {Factures et statut de reglement}
  {fig:s5-factures}

L'etablissement d'une facture, figure~\ref{fig:s5-creation}, ne part pas
d'une page blanche. L'ecran propose les \essentiel{missions facturables}
du client retenu, c'est-a-dire celles qui sont terminees et qui ne sont
rattachees a aucune facture. Le montant est preleve sur le devis accepte
au sprint~3. Cette continuite entre le devis et la facture est ce qui
supprime la double saisie que l'entreprise pratiquait auparavant entre son
tableur de suivi et son logiciel comptable, et avec elle les ecarts qu'une
recopie manuelle finit toujours par produire.

\capture{s5-creation-facture.jpg}
  {Etablissement d'une facture a partir des missions facturables}
  {fig:s5-creation}

\subsection{La surveillance}

Le centre d'alertes, figure~\ref{fig:s5-alertes}, rassemble en un seul
point ce qui reclame une decision : une assurance qui expire, une visite
technique depassee, un vehicule dont la revision est due, un equipement
au-dela de son seuil d'usure. Chaque alerte porte une priorite et un
statut, et peut etre rattachee au vehicule, au chauffeur, a l'equipement
ou a la mission concernes. Le filtre par priorite est ici l'element
determinant : sans lui, un centre d'alertes se remplit et cesse d'etre
consulte, ce qui le rend plus nuisible qu'utile.

\capture{s5-centre-alertes.jpg}
  {Centre d'alertes metier, filtrable par priorite et par statut}
  {fig:s5-alertes}

Le journal metier, figure~\ref{fig:s5-iboa}, releve d'une logique
differente. Il ne signale pas ce qu'il faut faire : il conserve ce qui
s'est passe. Chaque evenement publie par un service sur le courtier de
messages y est consigne avec sa criticite et son horodatage. L'ecran sert
a l'analyse apres coup --- comprendre pourquoi une mission a derape,
retrouver la chronologie d'un incident --- et non au pilotage courant.
C'est cette difference de finalite qui a justifie deux modules distincts
plutot qu'un seul.

\capture{s5-journal-iboa.jpg}
  {Journal des evenements metier et de leur criticite}
  {fig:s5-iboa}

% =====================================================================
\section{Tests de validation}

\begin{table}[H]
\centering
\caption{Tests de validation du sprint 5}
\label{tab:tests-s5}
\small
\begin{tabularx}{\textwidth}{@{}cXp{2.9cm}c@{}}
\toprule
\textbf{N\textsuperscript{o}} & \textbf{Cas teste} & \textbf{Attendu} &
\textbf{Etat} \\ \midrule
T1 & Mission deja facturee & Absente de la liste des facturables & Reussi \\
T2 & Facture etablie & Montants hors taxe, taxe et TTC coherents & Reussi \\
T3 & Echeance depassee sans reglement & Facture signalee en retard & Reussi \\
T4 & Controles metier lances deux fois & Aucune alerte en double & Reussi \\
T5 & Evenement publie sur le courtier & Consigne au journal avec sa criticite & Reussi \\
T6 & Alerte critique produite & Diffusion hors application declenchee & Reussi \\
\bottomrule
\end{tabularx}
\end{table}

% =====================================================================
\section*{Conclusion}
\addcontentsline{toc}{section}{Conclusion}

Ce sprint livre les trois services d'aval : facturation sans double
emploi possible, surveillance automatique des echeances et des seuils, et
journal metier alimente par le courtier de messages. Les quarante-sept
points du carnet ont ete realises.

Il met aussi en lumiere une \essentiel{limite du systeme d'evenements}.
La bibliotheque commune declare vingt-deux evenements d'integration ;
dans l'etat actuel, \essentiel{un seul service en publie et un seul en
consomme}. Le reste est declare sans emetteur ni destinataire. Ce n'est
pas un defaut de fonctionnement, mais une reserve de couplage qui n'a pas
ete exploitee faute de temps. Le chapitre de conclusion y revient au titre
des perspectives.

Le dernier sprint aborde ce qui reste : la collaboration entre les
utilisateurs et les deux services d'aide a la decision non encore
traites.

% % =====================================================================
%  Chapitre 9 - Sprint 6 : collaboration et aide a la decision
% =====================================================================
\chapter{Sprint 6 : collaboration et aide a la decision}
\label{chap:sprint6}

\section*{Introduction}
\addcontentsline{toc}{section}{Introduction}

Ce dernier sprint traite de trois services qui ne relevent d'aucun
domaine metier au sens strict, mais qui conditionnent l'usage quotidien
de la plateforme : la messagerie interne, l'assistant documentaire et le
moteur d'optimisation de flotte. Les deux derniers completent l'ensemble
des quatre services d'aide a la decision, dont l'estimation tarifaire et
la prediction de pannes ont deja ete traitees.

% =====================================================================
\section{Planification du sprint}

\subsection{Objectifs}

\begin{enumerate}
  \item Doter la plateforme d'une messagerie interne en temps reel.
  \item Offrir au client un assistant capable de repondre dans sa langue.
  \item Assister le gestionnaire dans l'affectation des vehicules et des
    chauffeurs.
  \item Consolider les indicateurs d'exploitation dans un tableau de
    bord.
\end{enumerate}

\subsection{Carnet du sprint}

\begin{table}[H]
\centering
\caption{Carnet du sprint 6}
\label{tab:backlog-s6}
\small
\begin{tabularx}{\textwidth}{@{}cp{2.3cm}Xcc@{}}
\toprule
\textbf{N\textsuperscript{o}} & \textbf{Acteur} & \textbf{Besoin exprime} &
\textbf{Prio.} & \textbf{Pts} \\ \midrule
1 & Utilisateur & Echanger en conversation privee ou en groupe & Should & 8 \\
2 & Utilisateur & Joindre un fichier a un message & Should & 5 \\
3 & Utilisateur & Voir qui est connecte & Could & 3 \\
4 & Client & Interroger un assistant sur les services & Could & 8 \\
5 & Client & Obtenir une reponse dans sa propre langue & Could & 5 \\
6 & Administrateur & Reindexer la base documentaire & Could & 3 \\
7 & Responsable & Obtenir une proposition d'affectation & Could & 8 \\
8 & Responsable & Consulter le tableau de bord consolide & Should & 5 \\
\midrule
\multicolumn{4}{@{}r}{\textbf{Total}} & \textbf{45} \\
\bottomrule
\end{tabularx}
\end{table}

% =====================================================================
\section{Messagerie interne}

\figReelle{chat-usecase.png}
  {Cas d'utilisation du service de messagerie}
  {fig:uc-chat}{1.00}

\figReelle{chat-sequence.png}
  {Sequence : envoi d'un message en temps reel}
  {fig:seq-chat}{1.00}

La messagerie est le seul service de la plateforme a maintenir une
\essentiel{liaison permanente} avec le navigateur. La
figure~\ref{fig:seq-chat} en montre la consequence : un message envoye
est d'abord persiste, puis diffuse aux participants connectes sans que
leur navigateur ait a interroger le serveur.

Ce choix a une contrepartie sur l'infrastructure. La passerelle doit
\essentiel{laisser passer la negociation de protocole} plutot que la
traiter comme une requete ordinaire, faute de quoi la liaison retombe sur
un mode degrade. C'est le seul endroit du systeme ou la configuration du
proxy porte une regle specifique.

\figReelle{chat-class-domain.png}
  {Modele du domaine du service de messagerie}
  {fig:dom-chat}{0.60}

Le modele repose sur trois classes. \exemple{Conversation} porte le type,
privee ou de groupe. \exemple{Participant} est une classe d'association
entre la conversation et l'utilisateur : elle porte la date de derniere
lecture, l'epinglage et l'activation des notifications, c'est-a-dire
\essentiel{ce qui est propre a chaque membre} et non a la conversation.
\exemple{Message} porte le contenu, un eventuel fichier joint et les
marqueurs de lecture.

% =====================================================================
\section{Assistant documentaire}

\figReelle{ayshaa-rag-usecase.png}
  {Cas d'utilisation de l'assistant documentaire}
  {fig:uc-rag}{1.00}

\figReelle{ayshaa-rag-sequence.png}
  {Sequence : traitement d'une question posee a l'assistant}
  {fig:seq-rag}{1.00}

L'assistant repond a partir d'une base documentaire vectorisee. La
sequence de la figure~\ref{fig:seq-rag} suit quatre etapes : detection de
la langue, recherche dans la base, consultation eventuelle des donnees
metier, puis redaction de la reponse.

\paragraph{Une recherche hybride.}
La recherche combine deux methodes. La premiere compare des
\essentiel{vecteurs de sens} : elle retrouve un document dont le
vocabulaire differe de celui de la question. La seconde est
\essentiel{lexicale} : elle privilegie la correspondance exacte des
termes. Les deux se completent, la premiere rattrapant les reformulations
et la seconde les references precises comme un numero de facture.

\paragraph{Trois langues, dont le dialecte local.}
Le detecteur de langue distingue le francais, l'anglais et l'arabe, en
prenant en charge \essentiel{le dialecte tunisien et sa transcription en
caracteres latins}. Ce point a justifie l'ecriture d'un detecteur propre
au projet : les detecteurs generiques traitent mal des messages courts et
melangeant les langues, ce qui est la norme dans les echanges
commerciaux locaux.

\paragraph{Un repli assume.}
La redaction finale s'appuie sur un modele de generation lorsqu'il est
disponible. \essentiel{S'il ne l'est pas, l'assistant ne tombe pas en
panne} : il restitue les extraits documentaires les plus pertinents. La
reponse est moins fluide, mais elle reste exacte et sourcee.

% =====================================================================
\section{Optimisation de l'affectation}

\figReelle{ai-optimizer-usecase.png}
  {Cas d'utilisation du moteur d'optimisation}
  {fig:uc-optimizer}{1.00}

\figReelle{ai-optimizer-sequence.png}
  {Sequence : production d'un plan d'affectation}
  {fig:seq-optimizer}{1.00}

Le moteur d'optimisation propose une affectation de vehicules et de
chauffeurs aux missions en attente. Il procede en deux temps
\essentiel{qu'il importe de distinguer} : d'abord il ecarte les
candidatures impossibles, ensuite il note celles qui restent.

\begin{table}[H]
\centering
\caption{Contraintes eliminatoires et criteres de notation}
\label{tab:optimisation}
\small
\begin{tabularx}{\textwidth}{@{}p{3.6cm}X@{}}
\toprule
\textbf{Nature} & \textbf{Regles appliquees} \\ \midrule
Contraintes eliminatoires & Classe du vehicule adaptee au chargement,
  categorie de permis du chauffeur, vehicule non immobilise pour
  maintenance, absence de chevauchement d'horaires, amplitude de conduite
  respectee \\
Criteres de notation & Proximite, disponibilite, adequation du vehicule a
  la mission, equilibrage de la charge entre chauffeurs \\
\bottomrule
\end{tabularx}
\end{table}

\essentiel{Le moteur propose, il ne decide pas.} Le plan produit est
soumis au gestionnaire, qui l'applique ou le modifie. Ce parti pris est le
meme que celui retenu pour l'estimation tarifaire au sprint~3 : l'aide a
la decision reste une aide.

L'echange est \essentiel{asynchrone} : la demande d'optimisation est
publiee sur le courtier de messages, le moteur la consomme et republie le
plan. Le gestionnaire n'attend pas devant son ecran, ce qui est acceptable
pour une operation de planification et impossible pour une estimation de
prix.

% =====================================================================
\section{Realisation}

\subsection{Les outils de collaboration}

La messagerie interne, figure~\ref{fig:s6-messagerie}, remplace les
groupes de discussion externes que l'entreprise utilisait pour coordonner
ses equipes. Le gain n'est pas le confort mais le \essentiel{rattachement
au contexte metier} : une conversation est liee a une mission, a un
vehicule ou a un ordre de travail, si bien que l'echange reste retrouvable
depuis l'objet qu'il concerne. Sur un fil de discussion generaliste, la
meme information se perd en quelques jours sous le flux des messages
suivants.

\capture{s6-messagerie.jpg}
  {Messagerie interne rattachee au contexte metier}
  {fig:s6-messagerie}

L'assistant documentaire du portail client, figure~\ref{fig:s6-assistant},
repond aux questions courantes sur les services proposes, les zones
desservies et les modalites de reservation. Il s'appuie sur les documents
de l'entreprise et \essentiel{cite ses sources}, ce qui permet au visiteur
de verifier la reponse. Son interet pour l'exploitation est de filtrer les
sollicitations les plus repetitives, qui occupaient jusque-la une part
non negligeable du temps telephonique.

\capture{s6-assistant-ayshaa.jpg}
  {Assistant documentaire du portail client}
  {fig:s6-assistant}

\subsection{Les outils d'aide a la decision}

Le plan d'affectation, figure~\ref{fig:s6-optimisateur}, restitue le
travail du moteur decrit a la section precedente. Chaque mission en
attente y recoit une proposition de vehicule et de chauffeur, assortie
de la note obtenue et des criteres qui l'ont determinee. Le gestionnaire
conserve la main : il applique le plan tel quel, ou il substitue une
affectation. Cette reversibilite est ce qui rend l'outil acceptable en
exploitation, ou une contrainte connue du seul exploitant --- un client
difficile, un chauffeur en formation --- echappe par nature au moteur.

\capture{s6-optimisateur.jpg}
  {Plan d'affectation propose par le moteur d'optimisation}
  {fig:s6-optimisateur}

Le tableau de bord, figure~\ref{fig:s6-bord1}, consolide en une page les
indicateurs d'exploitation : missions en cours et terminees, taux
d'utilisation du parc, vehicules immobilises, alertes ouvertes. Il agrege
des donnees produites par plusieurs services, ce qui en fait le point de
la plateforme ou la separation en services independants est la moins
visible pour l'utilisateur --- et c'est bien le resultat recherche.

\capture{s6-tableau-bord-1.jpg}
  {Tableau de bord : indicateurs d'exploitation consolides}
  {fig:s6-bord1}

% =====================================================================
\section{Tests de validation}

\begin{table}[H]
\centering
\caption{Tests de validation du sprint 6}
\label{tab:tests-s6}
\small
\begin{tabularx}{\textwidth}{@{}cXp{2.9cm}c@{}}
\toprule
\textbf{N\textsuperscript{o}} & \textbf{Cas teste} & \textbf{Attendu} &
\textbf{Etat} \\ \midrule
T1 & Message envoye a un groupe & Recu sans rechargement de page & Reussi \\
T2 & Fichier joint a un message & Telechargeable par les participants & Reussi \\
T3 & Question posee en dialecte tunisien & Reponse dans la meme langue & Reussi \\
T4 & Modele de generation indisponible & Extraits documentaires restitues & Reussi \\
T5 & Chauffeur sans le permis requis & Ecarte du plan d'affectation & Reussi \\
T6 & Vehicule en maintenance & Ecarte du plan d'affectation & Reussi \\
\bottomrule
\end{tabularx}
\end{table}

% =====================================================================
\section*{Conclusion}
\addcontentsline{toc}{section}{Conclusion}

Ce dernier sprint livre la messagerie interne, l'assistant documentaire
et le moteur d'optimisation. Les quarante-cinq points du carnet ont ete
realises, ce qui porte a trois cent six le total des points du projet.

Une ligne de conduite commune relie les quatre services d'aide a la
decision developpes au fil des sprints : \essentiel{aucun ne decide a la
place de l'utilisateur}. L'estimation propose un prix que le gestionnaire
peut modifier, la prediction signale un risque en le justifiant, l'assistant
cite ses sources, l'optimisation soumet un plan. Cette retenue n'est pas
une limite technique mais un choix : dans une entreprise de transport, la
responsabilite d'immobiliser un vehicule ou d'engager un tarif reste
humaine.

La realisation etant achevee, la conclusion generale dresse le bilan du
projet et en ouvre les perspectives.

% % =====================================================================
%  Conclusion generale
% =====================================================================
\cleardoublepage
\phantomsection
\addcontentsline{toc}{chapter}{Conclusion generale}
\chapter*{Conclusion generale}
\markboth{Conclusion generale}{}

\section*{Rappel de la demarche}
\addcontentsline{toc}{section}{Rappel de la demarche}

Ce projet de fin d'etudes partait d'une question posee en introduction :
comment doter une entreprise de transport de marchandises d'un systeme
d'information unifie qui centralise l'exploitation de sa flotte,
industrialise sa maintenance, ouvre un canal numerique a ses clients et
mette l'aide a la decision au service de ses choix quotidiens ?

L'etude prealable a montre que les pratiques en place, feuilles de
calcul, echanges informels et tableau de bord alimente a la main,
atteignaient leur limite des que plusieurs personnes devaient consulter
ou modifier la meme information. Aucune des trois familles d'outils
examinees ne couvrait l'ensemble du besoin, les manques les plus nets
portant sur l'aide a la decision et sur la surveillance des echeances
reglementaires.

L'analyse a traduit ce constat en exigences, puis le travail a ete
conduit selon le cadre Scrum, en six sprints d'un mois environ, chacun
livrant un domaine metier utilisable de bout en bout.

\section*{Resultats obtenus}
\addcontentsline{toc}{section}{Resultats obtenus}

Le perimetre realise comprend \essentiel{huit services metier} en .NET~8
appliquant la Clean Architecture, \essentiel{quatre services d'aide a la
decision} en Python, \essentiel{trois applications} destinees a des
publics distincts, et une infrastructure de vingt conteneurs orchestres
par Docker Compose autour d'une base relationnelle, d'un courtier de
messages et d'une base vectorielle. Les six sprints totalisent
\essentiel{trois cent six points de carnet, integralement livres}.

Les quatre objectifs fixes au depart ont ete atteints. Le referentiel est
\essentiel{unifie} : vehicules, equipements, personnels et clients sont
decrits une seule fois. La maintenance est \essentiel{industrialisee} :
ordres de travail suivis par etat, interventions tracees, stock surveille.
Le client dispose d'un \essentiel{canal autonome}, de la demande au suivi
en direct. Enfin, l'aide a la decision est \essentiel{integree au flux de
travail} plutot que releguee dans un module separe.

\section*{Difficultes rencontrees}
\addcontentsline{toc}{section}{Difficultes rencontrees}

Trois difficultes ont marque le deroulement du projet.

\paragraph{Le decoupage en services.}
Determiner ou passe la frontiere entre deux services est plus delicat que
de la coder. La question s'est posee des le sprint~3 : la demande de
transport releve-t-elle du domaine des missions ou de celui de la
facturation ? La reponse retenue, la rattacher aux missions, decoule de
l'observation que son cycle de vie s'acheve a la creation de la mission,
et non au reglement.

\paragraph{La coherence du vocabulaire entre services.}
Chaque service definit ses propres enumerations. L'un d'eux, le service
d'alertes, a fait diverger la sienne de celle du projet commun. Ecrire une
valeur du vocabulaire commun qui n'y figurait pas suffisait a faire
echouer la lecture. \essentiel{Cette difficulte est structurelle} : elle
est le revers de l'autonomie des services.

\paragraph{La restitution graphique de la modelisation.}
Produire des diagrammes UML lisibles s'est revele un travail a part
entiere. Certains atteignaient onze mille pixels de large, d'autres
presentaient des associations contournant la frontiere du systeme. Il a
fallu etablir des regles de mise en page explicites, puis les appliquer
de maniere systematique, pour que les modeles soient exploitables dans un
document imprime.

\section*{Apports personnels}
\addcontentsline{toc}{section}{Apports personnels}

Sur le plan \essentiel{technique}, ce projet a permis de mettre en oeuvre
une architecture microservices complete, de la conception au deploiement
conteneurise, avec deux piles technologiques et trois types
d'applications clientes. La Clean Architecture, appliquee huit fois, a
cesse d'etre une notion pour devenir une habitude de travail.

Sur le plan \essentiel{methodologique}, la conduite en sprints a montre
sa valeur au moment ou le perimetre a evolue : plusieurs fonctionnalites,
dont l'archivage reversible des comptes et la maintenance predictive, ont
ete identifiees en cours de projet. Un cycle sequentiel les aurait
renvoyees a une version ulterieure.

Sur le plan \essentiel{professionnel}, l'exercice le plus formateur a ete
d'apprendre a \essentiel{documenter ce qui ne fonctionne pas comme prevu}
plutot qu'a le passer sous silence. Les ecarts par rapport au patron de
reference, base partagee, absence de cache et de chiffrement, sont
consignes dans les diagrammes eux-memes. C'est une pratique que
l'entreprise d'accueil, \entreprise, encourage, et qui rend une
architecture reprenable par un tiers.

\section*{Perspectives}
\addcontentsline{toc}{section}{Perspectives}

Le travail realise ouvre plusieurs prolongements, classes ici par ordre
de priorite.

\begin{table}[H]
\centering
\caption{Perspectives d'evolution}
\label{tab:perspectives}
\small
\begin{tabularx}{\textwidth}{@{}p{3.4cm}X@{}}
\toprule
\textbf{Axe} & \textbf{Contenu} \\ \midrule
Mise en securite & Terminaison du chiffrement a la passerelle et
  redirection du trafic en clair : prealable a toute mise en ligne \\
Exploitation du bus & Sur vingt-deux evenements declares, un seul est
  publie et un seul consomme ; les exploiter reduirait les lectures
  directes entre domaines \\
Cloisonnement des donnees & Passer a une base par service, en remplacant
  les lectures en SQL direct par des appels ou des evenements \\
Bascule vers l'apprentissage & Substituer un modele appris au moteur de
  regles des qu'un historique suffisant sera constitue \\
Mise a l'echelle & Cache applicatif et orchestrateur de grappe \\
Application chauffeur & Mise en file locale des positions lorsque la
  couverture reseau fait defaut \\
\bottomrule
\end{tabularx}
\end{table}

\section*{Mot de la fin}
\addcontentsline{toc}{section}{Mot de la fin}

La plateforme livree n'est pas un prototype : elle fonctionne, elle est
deployee et elle repond aux besoins exprimes. Elle n'est pas non plus un
produit acheve, et le present rapport s'est attache a dire ou passe cette
limite. Cette distinction est peut-etre l'apprentissage principal de ces
\stageDuree{} mois : un systeme d'information se juge autant a ce qu'il
assume ne pas faire qu'a ce qu'il realise.


% =====================================================================
%  BIBLIOGRAPHIE ET ANNEXES
% =====================================================================
% % =====================================================================
%  Bibliographie et netographie
% =====================================================================
%  Tenue a la main plutot qu'avec biblatex : le guide impose un format
%  precis (auteurs, titre, editeur, date, URL et date de consultation) et
%  la compilation reste un simple pdflatex lance deux fois.
% =====================================================================
\cleardoublepage
\phantomsection
\addcontentsline{toc}{chapter}{Bibliographie et netographie}
\chapter*{Bibliographie et netographie}
\markboth{Bibliographie et netographie}{}

\section*{Ouvrages}
\addcontentsline{toc}{section}{Ouvrages}

\begin{enumerate}[label={[\arabic*]},leftmargin=*,itemsep=4pt]

\item MARTIN Robert C., \textit{Clean Architecture: A Craftsman's Guide
  to Software Structure and Design}, Prentice Hall, 2017, 432 p.

\item NEWMAN Sam, \textit{Building Microservices: Designing Fine-Grained
  Systems}, 2\textsuperscript{e} edition, O'Reilly Media, 2021, 616 p.

\item EVANS Eric, \textit{Domain-Driven Design: Tackling Complexity in
  the Heart of Software}, Addison-Wesley, 2003, 560 p.

\item ROQUES Pascal, \textit{UML 2 par la pratique : etudes de cas et
  exercices corriges}, 7\textsuperscript{e} edition, Eyrolles, 2009,
  396 p.

\item AUBRY Claude, \textit{Scrum : le guide pratique de la methode agile
  la plus populaire}, 5\textsuperscript{e} edition, Dunod, 2019, 336 p.

\end{enumerate}

\section*{Articles et specifications}
\addcontentsline{toc}{section}{Articles et specifications}

\begin{enumerate}[label={[\arabic*]},leftmargin=*,itemsep=4pt,start=6]

\item CHEN Tianqi et GUESTRIN Carlos, \guillemotleft{} XGBoost: A Scalable
  Tree Boosting System \guillemotright{}, \textit{Proceedings of the
  22\textsuperscript{nd} ACM SIGKDD International Conference on Knowledge
  Discovery and Data Mining}, San Francisco, 2016, p. 785-794.

\item LEWIS Patrick \textit{et al.}, \guillemotleft{} Retrieval-Augmented
  Generation for Knowledge-Intensive NLP Tasks \guillemotright{},
  \textit{Advances in Neural Information Processing Systems}, vol. 33,
  2020, p. 9459-9474.

\item OBJECT MANAGEMENT GROUP, \textit{Unified Modeling Language
  Specification, version 2.5.1}, decembre 2017.

\item SCHWABER Ken et SUTHERLAND Jeff, \textit{The Scrum Guide}, novembre
  2020, 13 p.

\end{enumerate}

\section*{Netographie}
\addcontentsline{toc}{section}{Netographie}

\begin{enumerate}[label={[\arabic*]},leftmargin=*,itemsep=4pt,start=10]

\item MICROSOFT, \textit{Documentation .NET 8 et ASP.NET Core},
  \url{https://learn.microsoft.com/dotnet/}, consulte en mars 2026.

\item MICROSOFT, \textit{Entity Framework Core},
  \url{https://learn.microsoft.com/ef/core/}, consulte en mars 2026.

\item MASSTRANSIT, \textit{Documentation officielle},
  \url{https://masstransit.io/documentation}, consulte en avril 2026.

\item META, \textit{Documentation React},
  \url{https://react.dev/}, consulte en avril 2026.

\item VERCEL, \textit{Documentation Next.js},
  \url{https://nextjs.org/docs}, consulte en mai 2026.

\item EXPO, \textit{Documentation Expo et React Native},
  \url{https://docs.expo.dev/}, consulte en mai 2026.

\item POSTGRESQL GLOBAL DEVELOPMENT GROUP, \textit{PostgreSQL 16
  Documentation}, \url{https://www.postgresql.org/docs/16/}, consulte en
  mars 2026.

\item BROADCOM, \textit{RabbitMQ Documentation},
  \url{https://www.rabbitmq.com/docs}, consulte en avril 2026.

\item QDRANT, \textit{Qdrant Documentation},
  \url{https://qdrant.tech/documentation/}, consulte en juillet 2026.

\item DOCKER INC., \textit{Docker Compose Documentation},
  \url{https://docs.docker.com/compose/}, consulte en mars 2026.

\item PLANTUML, \textit{Reference Guide},
  \url{https://plantuml.com/guide}, consulte en juin 2026.

\end{enumerate}

% \appendix
% \input{annexes/A-extraits-code}
% \input{annexes/B-journal-de-bord}
% \input{annexes/C-jeux-de-tests}
% \input{annexes/D-diagrammes-complementaires}
% % =====================================================================
%  Resume et abstract, en quatrieme de couverture
% =====================================================================
\cleardoublepage
\thispagestyle{empty}
\vspace*{1cm}

\begin{center}
  {\Large\sffamily\bfseries\color{bleuTitre} Resume\par}
\end{center}
\vspace{0.5cm}

Le transport routier de marchandises repose sur la coordination de
quatre ressources aux cycles de vie independants : un parc de vehicules,
des equipements embarques, des chauffeurs et un portefeuille de clients.
Les entreprises de taille moyenne gerent cette coordination avec des
outils qui n'y sont pas destines, ce qui produit une information
dispersee, une maintenance subie et une relation client opaque.

Ce projet de fin d'etudes, realise au sein de \entreprise, a consiste a
concevoir et developper \sujetCourt, une plateforme de gestion de flotte
et de transport de marchandises fondee sur une architecture
microservices. Elle reunit huit services metier en .NET~8 appliquant la
Clean Architecture, quatre services d'aide a la decision en Python, et
trois applications destinees a des publics distincts : un back-office, un
portail client et une application mobile pour les chauffeurs. Le tout est
deploye en vingt conteneurs orchestres par Docker Compose.

Le travail a ete conduit selon le cadre Scrum, en six sprints couvrant
successivement les identites et les acces, le referentiel de flotte, la
chaine commerciale de la demande au suivi en direct, la maintenance
curative et predictive, la facturation et les alertes, enfin la
collaboration et l'aide a la decision.

\vspace{0.4cm}
\noindent
\textbf{Mots-cles :} microservices, gestion de flotte, transport de
marchandises, Clean Architecture, Scrum, aide a la decision.

\vspace{1.4cm}
{\color{grisFilet}\rule{\textwidth}{0.8pt}}
\vspace{1.4cm}

\begin{center}
  {\Large\sffamily\bfseries\color{bleuTitre} Abstract\par}
\end{center}
\vspace{0.5cm}

Road freight transport relies on coordinating four resources with
independent life cycles: a vehicle fleet, on-board equipment, drivers and
a customer portfolio. Medium-sized companies handle this coordination
with tools that were never designed for it, which leads to scattered
information, reactive maintenance and limited visibility for customers.

This final-year project, carried out at \entreprise, consisted in
designing and developing \sujetCourt, a fleet and freight transport
management platform built on a microservices architecture. It brings
together eight .NET~8 business services following Clean Architecture,
four Python decision-support services, and three applications serving
distinct audiences: a back-office, a customer portal and a mobile
application for drivers. The whole system is deployed as twenty
containers orchestrated by Docker Compose.

The work followed the Scrum framework, over six sprints covering
identity and access management, the fleet repository, the commercial
chain from request to live tracking, corrective and predictive
maintenance, invoicing and business alerts, and finally collaboration and
decision support.

\vspace{0.4cm}
\noindent
\textbf{Keywords:} microservices, fleet management, freight transport,
Clean Architecture, Scrum, decision support.


\end{document}
